%% =====================================================================
%% The packages and settings the template itself needs, in the order they must be loaded (main.tex reads
%% this file first, before your settings/packages.tex). Your own packages go in settings/packages.tex;
%% change this file only to change the template.
%% =====================================================================

%% Maths and units
\ifthesisotfmath\else\usepackage{bm}\fi % https://ctan.org/pkg/bm %% bold symbols (\vec, \tensor); the OpenType maths fonts have their own \bm (MathTime also \hm, heavy)
\usepackage{siunitx}            % https://ctan.org/pkg/siunitx %% quantities and units: \qty{}{}, \num{}, \unit{}
\usepackage{derivative}         % https://ctan.org/pkg/derivative %% derivatives and differentials: \odv{y}{x}, \pdv{u}{t}, \odif{x} (see commands.tex)
%% =====================================================================
%% Crossing out terms in equations: \cancel{a}, \bcancel{a}, \xcancel{a}, \cancelto{0}{x}
%% (cancel package), made fit for a tagged PDF:
%% - the lines are drawn by pict2e as lines, not with the LaTeX line font that cancel uses
%%   (its characters have no Unicode meaning, which PDF/UA does not allow); the look is the same;
%% - the MathML says what is crossed out: <menclose notation="updiagonalstrike"> around the term
%%   (without this the term would be lost for screen readers), and for \cancelto an arrow with
%%   the value as a superscript.
%% =====================================================================
\usepackage{pict2e}             % https://ctan.org/pkg/pict2e
\usepackage{cancel}             % https://ctan.org/pkg/cancel
\makeatletter
\let\OriginalPictureCmds\relax  % cancel's lines from pict2e (pict2e loaded first would make cancel use the font)

%% The crossed-out term (and the \cancelto value) is saved for the MathML at the end of the small
%% formula in which cancel sets it.
\patchcmd\@cancel{$#2{#3}$}{$#2{#3}\thesis@cancelsave#2$}{}
    {\PackageWarning{thesis}{Could not patch \string\cancel}}
\patchcmd\@cancelto{$#3{#4}$}{$#3{#4}\thesis@cancelsave#3$}{}
    {\PackageWarning{thesis}{Could not patch \string\cancelto}}
\patchcmd\canto@vector{$#7{#8}\m@th$}{$#7{#8}\thesis@canceltosave#7\m@th$}{}
    {\PackageWarning{thesis}{Could not patch \string\cancelto}}

\ExplSyntaxOn
%% Without MathML (untagged editions) nothing is saved or added.
\cs_new_protected:Npn \__thesis_cancel_save:nN #1#2 { }
\cs_new_protected:Npn \__thesis_cancel_mathml:nn #1#2 { #2 }
\AddToHook{begindocument}
  {
    \IfPackageLoadedT{luamml}
      {
        %% save the formula's MathML under a label (not in text mode, where cancel uses \hbox); it is
        %% only a part of the MathML of the formula around it, so it gets no structure of its own
        \cs_set_protected:Npn \__thesis_cancel_save:nN #1#2
          {
            \token_if_eq_meaning:NNF #2 \hbox
              {
                \luamml_save:nNn {#1} #2 { mrow }
                \thesis_luamml_no_structure:
              }
          }
        %% #1: the MathML (a Lua table) that replaces what #2 sets
        \cs_set_protected:Npn \__thesis_cancel_mathml:nn #1#2
          { \luamml_annotate:nen { 1 } { nucleus = true, mathml = {#1} } {#2} }
      }
  }
\cs_new_protected:Npn \thesis@cancelsave { \__thesis_cancel_save:nN { thesiscancel } }
\cs_new_protected:Npn \thesis@canceltosave { \__thesis_cancel_save:nN { thesiscancelto } }

%% #1 MathML notation, #2 cancel's lines, #3 maths style, #4 the term
\cs_new_protected:Npn \__thesis_cancel:nnNn #1#2#3#4
  {
    \__thesis_cancel_mathml:nn { [0] = 'menclose', notation = '#1', consume_label'thesiscancel' }
      { \@cancel {#2} #3 {#4} }
  }
\DeclareRobustCommand \cancel [1]
  {
    \mode_if_math:TF
      { \mathpalette { \__thesis_cancel:nnNn { updiagonalstrike } { \@can@slash { } } } {#1} }
      { \@cancel { \@can@slash { } } \hbox {#1} }
  }
\DeclareRobustCommand \bcancel [1]
  {
    \mode_if_math:TF
      { \mathpalette { \__thesis_cancel:nnNn { downdiagonalstrike } { \@can@slash { - } } } {#1} }
      { \@cancel { \@can@slash { - } } \hbox {#1} }
  }
\DeclareRobustCommand \xcancel [1]
  {
    \mode_if_math:TF
      {
        \mathpalette
          { \__thesis_cancel:nnNn { updiagonalstrike~downdiagonalstrike } { \@can@slash { + } \@can@slash { - } } }
          {#1}
      }
      { \@cancel { \@can@slash { + } \@can@slash { - } } \hbox {#1} }
  }
%% \cancelto{value}{term}: #1 style of the value, #2 value, #3 style of the term, #4 term
\cs_new_protected:Npn \__thesis_cancelto:NnNn #1#2#3#4
  {
    \__thesis_cancel_mathml:nn
      {
        [0] = 'msup',
        { [0] = 'menclose', notation = 'updiagonalarrow', consume_label'thesiscancel' },
        consume_label'thesiscancelto'
      }
      { \@cancelto #1 {#2} #3 {#4} }
  }
\cs_set:Npn \cancelto #1#2 % the value one size smaller, as cancel does by default
  {
    \mathchoice
      { \__thesis_cancelto:NnNn \textstyle {#1} \displaystyle {#2} }
      { \__thesis_cancelto:NnNn \scriptstyle {#1} \textstyle {#2} }
      { \__thesis_cancelto:NnNn \scriptscriptstyle {#1} \scriptstyle {#2} }
      { \__thesis_cancelto:NnNn \scriptscriptstyle {#1} \scriptscriptstyle {#2} }
  }
\ExplSyntaxOff
\makeatother
         % crossing out terms: \cancel{a}, \bcancel, \xcancel, \cancelto{0}{x}
%% =====================================================================
%% Matrices with nicematrix: labels outside the matrix (first-row, last-col, ...), blocks
%% (\Block{2-2}{A}), braces over parts (\OverBrace, \UnderBrace, \SubMatrix after \CodeAfter),
%% frames, and equal column widths (NiceMatrixBlock with auto-columns-width).
%%   \begin{pNiceMatrix}[first-row,last-col]
%%     x & y & \\  1 & 0 & \text{row 1} \\  0 & 1 & \text{row 2}
%%   \end{pNiceMatrix}
%% Only the matrix environments (NiceMatrix, pNiceMatrix, bNiceMatrix, BNiceMatrix, vNiceMatrix,
%% VNiceMatrix) are made fit for a tagged PDF; for block matrices with lines use one of them
%% with \hline and the key vlines (e.g. [vlines=2]). Not NiceArray and NiceTabular, whose
%% content would be lost for screen readers; for tables use multirow and makecell.
%% amsmath's \cdots, \vdots, \ddots work in the cells as everywhere.
%%
%% Tagged PDF: nicematrix draws the matrix with PGF, so its MathML would be empty. The entries
%% are therefore set once more, out of sight, as an amsmath matrix, and that MathML is used,
%% made to match the matrix as seen (nicematrix.lua): labels and braces are read after the
%% matrix, blocks are cells spanning rows and columns.
%% =====================================================================
\usepackage{nicematrix}         % https://ctan.org/pkg/nicematrix
\makeatletter
\NewEnvironmentCopy{thesis@nm@NiceMatrix}{NiceMatrix}
\NewEnvironmentCopy{thesis@nm@pNiceMatrix}{pNiceMatrix}
\NewEnvironmentCopy{thesis@nm@bNiceMatrix}{bNiceMatrix}
\NewEnvironmentCopy{thesis@nm@BNiceMatrix}{BNiceMatrix}
\NewEnvironmentCopy{thesis@nm@vNiceMatrix}{vNiceMatrix}
\NewEnvironmentCopy{thesis@nm@VNiceMatrix}{VNiceMatrix}
\makeatother

\ExplSyntaxOn
\makeatletter
\bool_new:N \l__thesis_nm_mathml_bool
\AddToHook{begindocument}
  {
    \IfPackageLoadedT{luamml}
      {
        \bool_set_true:N \l__thesis_nm_mathml_bool
        \directlua{ dofile(kpse.find_file('template/nicematrix.lua')) }
      }
  }
\int_new:N \g__thesis_nm_int
\tl_new:N \l__thesis_nm_body_tl
\tl_new:N \l__thesis_nm_after_tl
\tl_new:N \l__thesis_nm_spec_tl
\box_new:N \l__thesis_nm_box

%% MathML of #2 (maths) under the label #1, for the MathML of the matrix only: no structure of its own
\cs_new_protected:Npn \__thesis_nm_save:nn #1#2
  {
    \hbox_set:Nn \l__thesis_nm_box
      {
        \__thesis_nm_no_structure:
        $ \m@th #2 \luamml_save:nNn {#1} \textstyle { mrow } \__thesis_nm_no_structure: $
      }
  }
\cs_new_protected:Npn \__thesis_nm_no_structure:
  {
    \int_set:Nn \l__luamml_flag_int
      { \l__luamml_flag_int - 8 * \int_mod:nn { \int_div_truncate:nn { \l__luamml_flag_int } { 8 } } { 2 } }
  }
%% every cell gets a marker, so that empty cells and rows keep their place in the MathML
\cs_new_protected:Npn \__thesis_nm_cell:
  { \luamml_annotate:nen { 1 } { nucleus = true, mathml = { [0] = 'mrow', ['thesis-cell'] = 'true' } } { { \hbox:n { } } } }
\cs_new:Npn \__thesis_nm_before_body:w #1 \Body #2 \q_stop { \exp_not:n {#2} }
\cs_new:Npn \__thesis_nm_split_after:w #1 \CodeAfter #2 \q_stop
  { \tl_set:Nn \l__thesis_nm_body_tl {#1} \tl_set:Nn \l__thesis_nm_after_tl {#2} }

%% a brace or sub-matrix after \CodeAfter: its kind, cells and label, for nicematrix.lua
\cs_new_protected:Npn \__thesis_nm_note:nnnn #1#2#3#4 % kind, from, to, label
  {
    \int_gincr:N \g__thesis_nm_int
    \tl_if_blank:nTF {#4}
      { \tl_put_right:Ne \l__thesis_nm_spec_tl { { kind='#1', from='\tl_to_str:n {#2}', to='\tl_to_str:n {#3}' }, } }
      {
        \__thesis_nm_save:nn { thesisnm \int_use:N \g__thesis_nm_int } {#4}
        \tl_put_right:Ne \l__thesis_nm_spec_tl
          {
            { kind='#1', from='\tl_to_str:n {#2}', to='\tl_to_str:n {#3}',
              label=consume_label'thesisnm \int_use:N \g__thesis_nm_int' },
          }
      }
  }
\cs_new_protected:Npn \__thesis_nm_key:nnn #1#2#3 % keys, key, Lua name
  { \clist_if_in:nnT {#1} {#2} { \tl_put_right:Nn \l__thesis_nm_spec_tl { #3 = true, } } }
\cs_new:Npn \__thesis_nm_delim:N #1
  { \str_case:nnF {#1} { {(} {p} {[} {b} {\{} {B} {|} {v} {\|} {V} {\vert} {v} {\Vert} {V} } {} }

%% The entries set once more as an amsmath matrix, saved as MathML under thesisnm.
%% #1: delimiter letter (p, b, B, v, V or empty), #2: keys, #3: body
\cs_new_protected:Npn \__thesis_nm_mathml:nnn #1#2#3
  {
    \group_begin:
      \tl_set:Nn \l__thesis_nm_body_tl {#3}
      \tl_clear:N \l__thesis_nm_after_tl
      \tl_if_in:NnT \l__thesis_nm_body_tl { \Body }
        { \tl_set:Ne \l__thesis_nm_body_tl { \exp_after:wN \__thesis_nm_before_body:w \l__thesis_nm_body_tl \q_stop } }
      \tl_if_in:NnT \l__thesis_nm_body_tl { \CodeAfter }
        { \exp_after:wN \__thesis_nm_split_after:w \l__thesis_nm_body_tl \q_stop }
      %% the keys for labels outside the matrix, and the language of the notes
      \tl_set:Ne \l__thesis_nm_spec_tl
        { lang = '\str_if_eq:eeTF { \languagename } { dutch } { nl } { en }', }
      \clist_map_inline:nn { {first-row}{first_row}, {last-row}{last_row}, {first-col}{first_col}, {last-col}{last_col} }
        { \__thesis_nm_key:nnn {#2} ##1 }
      \tl_put_right:Ne \l__thesis_nm_spec_tl { notes = \c_left_brace_str }
      %% braces and sub-matrices; everything else after \CodeAfter is drawing only
      \hbox_set:Nn \l__thesis_nm_box
        {
          \DeclareDocumentCommand \OverBrace { O{} m m m O{} } { \__thesis_nm_note:nnnn { over } {##2} {##3} {##4} }
          \DeclareDocumentCommand \UnderBrace { O{} m m m O{} } { \__thesis_nm_note:nnnn { under } {##2} {##3} {##4} }
          \DeclareDocumentCommand \SubMatrix { m m m m O{} }
            { \__thesis_nm_note:nnnn { \__thesis_nm_delim:N ##1 } {##2} {##3} { } }
          \DeclareDocumentCommand \line { O{} m m O{} } { }
          \DeclareDocumentCommand \tikz { O{} m } { }
          \DeclareDocumentEnvironment { tikzpicture } { O{} +b } { } { }
          \DeclareDocumentCommand \Block { O{} m D<>{} +m O{} } { }
          \l__thesis_nm_after_tl
          %% (the spec is filled inside this box: hand it out)
          \tl_gset_eq:NN \g_tmpa_tl \l__thesis_nm_spec_tl
        }
      \tl_set:Ne \l__thesis_nm_spec_tl { \g_tmpa_tl \c_right_brace_str }
      %% the matrix: nicematrix's commands give plain entries, a \Block a marker with its
      %% entry, its size and whether it is drawn
      \DeclareDocumentCommand \Block { O{} m D<>{} +m O{} }
        {
          \int_gincr:N \g__thesis_nm_int
          \tl_if_blank:nF {##4} { \__thesis_nm_save:nn { thesisnm \int_use:N \g__thesis_nm_int } {##4} }
          \tl_set:Nn \l_tmpa_tl { false }
          \tl_if_in:nnT {##1} { draw } { \tl_set:Nn \l_tmpa_tl { true } }
          \luamml_annotate:nen { 1 }
            {
              nucleus = true,
              mathml = { [0] = 'mrow', ['thesis-span'] = '\tl_to_str:n {##2}',
                ['thesis-draw'] = '\l_tmpa_tl'
                \tl_if_blank:nF {##4} { , consume_label'thesisnm \int_use:N \g__thesis_nm_int' } }
            }
            { { \hbox:n { } } }
        }
      \cs_set_eq:NN \Cdots \cdots \cs_set_eq:NN \Ldots \ldots \cs_set_eq:NN \Vdots \vdots
      \cs_set_eq:NN \Ddots \ddots \cs_set_eq:NN \Iddots \iddots
      \cs_set_eq:NN \Hline \hline \cs_set_eq:NN \hdottedline \hline
      \cs_set:Npn \NotEmpty { }
      \DeclareDocumentCommand \RowStyle { O{} m } { }
      \DeclareDocumentCommand \diagbox { m m } { ##1 }
      \cs_set:Npn \rotate { }
      \setcounter{MaxMatrixCols}{100}
      \tl_replace_all:Nnn \l__thesis_nm_body_tl { & } { \__thesis_nm_cell: & }
      \tl_replace_all:Nnn \l__thesis_nm_body_tl { \\ } { \__thesis_nm_cell: \\ }
      \tl_put_right:Nn \l__thesis_nm_body_tl { \__thesis_nm_cell: }
      \__thesis_nm_save:nn { thesisnm }
        { \exp_args:NnV \use:nn { \begin{#1matrix} } \l__thesis_nm_body_tl \end{#1matrix} }
      \tl_gset_eq:NN \g_tmpa_tl \l__thesis_nm_spec_tl
    \group_end:
    \tl_set_eq:NN \l__thesis_nm_spec_tl \g_tmpa_tl
  }

\clist_map_inline:nn { {}, p, b, B, v, V }
  {
    \RenewDocumentEnvironment { #1 NiceMatrix } { !O{} +b }
      {
        \bool_if:NTF \l__thesis_nm_mathml_bool
          {
            \__thesis_nm_mathml:nnn {#1} {##1} {##2}
            \luamml_annotate:nen { 1 }
              { nucleus = true, mathml = thesis_nicematrix(consume_label'thesisnm', { \l__thesis_nm_spec_tl }) }
              {
                {
                  \tagpdfsetup { table/tagging = false }
                  \int_set:Nn \l__luamml_flag_int { 0 } % (its own MathML is replaced)
                  \begin{thesis@nm@ #1 NiceMatrix}[##1] ##2 \end{thesis@nm@ #1 NiceMatrix}
                }
              }
          }
          { \begin{thesis@nm@ #1 NiceMatrix}[##1] ##2 \end{thesis@nm@ #1 NiceMatrix} }
      }
      { }
  }
\makeatother
\ExplSyntaxOff
     % matrices with labels outside, blocks, braces: pNiceMatrix etc.

%% Figures and tables
\usepackage{array}              % https://ctan.org/pkg/array %% extra column types for tables
\usepackage{booktabs}           % https://ctan.org/pkg/booktabs %% \toprule, \midrule, \bottomrule
\usepackage{xltabular}          % https://ctan.org/pkg/xltabular %% tabularx (wrapping columns) + longtable (across pages)
\usepackage{multirow}           % https://ctan.org/pkg/multirow %% table cells spanning several rows: \multirow{2}{*}{text}
\usepackage{makecell}           % https://ctan.org/pkg/makecell %% line breaks in a cell: \makecell{two\\lines}
\usepackage{placeins}           % https://ctan.org/pkg/placeins %% \FloatBarrier: floats placed so far go out before what follows
                                %  ([H], "here, no matter what", also works: float.sty comes with the algorithms below)
\usepackage{pdfpages}           % https://ctan.org/pkg/pdfpages %% whole PDF pages, e.g. a published paper: \includepdf[pages=-]{paper.pdf}
%% Sub-figures and sub-tables ((a), (b), ...) are part of the class: see thesis.cls.
%% =====================================================================
%% Tables: notes under a table, one text size, wrapping columns, sideways floats.
%% Read by settings/packages.tex. Part of the template: edit settings/packages.tex instead.
%% =====================================================================
%% Notes under a table, marked in the table with \tnote{a} (as with the threeparttable package, whose
%% notes break the tagged PDF). They take the width of the table just above them (a tabular), or the
%% width of the line (tabularx, xltabular), or the width given: \begin{tablenotes}[6cm].
%%   \begin{tablenotes} \item[a] Corrected for blockage. \item[b] ... \end{tablenotes}
\makeatletter
\newdimen\thesis@tabwd
\AddToHook{cmd/@floatboxreset/after}{\global\thesis@tabwd\z@}
\AddToHook{env/tabular/after}{\ifhmode\setbox\z@\lastbox\global\thesis@tabwd\wd\z@\box\z@\fi}
\AddToHook{env/tabularx/after}{\global\thesis@tabwd\linewidth}
\AddToHook{env/xltabular/after}{\global\thesis@tabwd\linewidth}
%% (a marker takes no width, so numbers in a column stay aligned and the marker sticks out to the right)
\providecommand*{\tnote}[1]{\rlap{\textsuperscript{#1}}}
\newenvironment{tablenotes}[1][\ifdim\thesis@tabwd>\z@\thesis@tabwd\else\linewidth\fi]
  {\par\vspace{0.4\onelineskip}\noindent\begin{minipage}[t]{#1}\footnotesize
   \begin{list}{}{\setlength{\leftmargin}{1em}\setlength{\labelwidth}{0.8em}\setlength{\labelsep}{0.2em}%
     \setlength{\itemsep}{\z@}\setlength{\parsep}{\z@}\setlength{\topsep}{\z@}\setlength{\partopsep}{\z@}%
     \let\makelabel\thesis@tnotelabel}}
  {\end{list}\end{minipage}\par}
\newcommand*{\thesis@tnotelabel}[1]{\textsuperscript{#1}\hfil}
\makeatother
%% All tables in one size, the caption size (\small): table text then sits between the captions and the text.
%% (The nomenclature lists are not tables in this sense and keep the text size.)
\setfloatadjustment{table}{\small} % applied inside every table float, also sideways ones (thesis.cls)
\AddToHook{env/xltabular/begin}{\small} % long tables are not floats
%% Wrapping (X) columns are set ragged right without hyphenation: in a narrow column justified text
%% leaves wide gaps, and a word broken over two lines of a table cell reads badly.
\renewcommand{\tabularxcolumn}[1]{>{\raggedright\arraybackslash}p{#1}}
%% Sideways floats (sidewaystable, sidewaysfigure) for material too wide for the page: a page of
%% their own, turned so that the top of the table is at the left edge on left and right pages alike
%% (the book convention: the reader always turns the book clockwise, and a table continued over a
%% spread reads as one page). The text keeps flowing around it.
\makeatletter
\usepackage[figuresright]{rotating} % https://ctan.org/pkg/rotating
%% Digital edition: those pages are shown turned upright in PDF viewers. A marker inside the float sets
%% /Rotate 90 on the page it lands on (only that page). Not in the print editions: for the printer, a rotated page can confuse the imposition software.
\let\thesisrotatepage\relax
\ifprintversion\else
    \ExplSyntaxOn
    \protected\def\thesisrotatepage{\pdfmanagement_add:nnn { ThisPage } { Rotate } { 90 }}
    \ExplSyntaxOff
    \AddToHook{env/sidewaystable/end}{\thesisrotatepage}
    \AddToHook{env/sidewaysfigure/end}{\thesisrotatepage}
\fi
\makeatother
       % notes under tables (\tnote, tablenotes), table text size, sideways floats
\usetikzlibrary{positioning,arrows.meta} % TikZ itself is loaded by the class (thumb tabs)

%% Text and code
\usepackage{enumitem}           % https://ctan.org/pkg/enumitem %% list settings: \begin{enumerate}[label=(\roman*)], [noitemsep]
\usepackage{orcidlink}          % https://ctan.org/pkg/orcidlink %% the ORCID iD icon, used by \orcid (commands.tex)
\usepackage{listings}           % https://ctan.org/pkg/listings %% code listings
%% =====================================================================
%% Tagged PDF (PDF/UA): what LaTeX's tagging code does not handle by itself (yet), and where its look
%% differs from the untagged editions. Read by template/packages.tex, right after listings is loaded; everything
%% here does something only when tagging is on (\DocumentMetadata{..., tagging=on}).
%%   1. Code listings
%%   2. Equation numbers in the MathML (with tagging.lua)
%%   3. Theorems and proofs straight after one another (LaTeX's code of June 2026)
%%   4. Figures side by side
%%   5. No "Rerun" warning for tagpdf's running totals
%%   6. Table cells: their scope and column span also on the cell itself (for PAC)
%% =====================================================================
\ExplSyntaxOn
\tag_if_active:T { \cs_gset:Npn \thesistagging {} }
\ExplSyntaxOff
\ifdefined\thesistagging
\makeatletter

%% ---------------------------------------------------------------------------------------------
%% 1. Code listings: a code block (a verbatim structure of code lines), the line numbers as labels and
%% the caption as the listing's caption. LaTeX brings this itself from November 2026 (latex-lab-listings,
%% loaded with listings); for the June 2026 release, the same code, from the LaTeX team's set-up for the
%% tagpdf documentation (tagpdfdocu-patches.sty, 2026-08-21, LaTeX Project Public License). Without it,
%% every listing closed the section around it in the tagged structure.
\IfPackageLoadedTF{latex-lab-listings}{\let\thesis@lablistings\relax}{}
\ifdefined\thesis@lablistings\else
\ExplSyntaxOn
\NewTaggingSocket{listings/number}{2}
\NewTaggingSocketPlug{listings/number}{default}
  {
    \tag_mc_end_push:
    \tagstructbegin{tag=Lbl}
    \tagmcbegin{}
    #2
    \tagmcend
    \tagstructend
    \tag_mc_begin_pop:n{}
  }
\AssignTaggingSocketPlug{listings/number}{default}
\ExplSyntaxOff
\DeclareInstance{blockenv}{lstlisting}{std}{%
    name                   = lstlisting,
    tag-name               = verbatim,
    tag-attr-class         = ,
    tagging-recipe         = standard,
    inner-level-counter    = ,
    transparent-level      = true,
    block-instance         = lstlisting,
    inner-instance         = ,
    final-code             = \AssignStructureRole{para/textblock}
                               {\UseStructureName{block/verbatim/codeline}}\ignorespaces,
    tagging-suppress-paras = true}
\tagpdfsetup{float/new=lstlisting}
\lst@Key{numbers}{none}{%
    \let\lst@PlaceNumber\@empty
    \lstKV@SwitchCases{#1}%
    {none:\\%
     left:\def\lst@PlaceNumber{\llap{\normalfont
                \UseTaggingSocket{listings/number}{}{\lst@numberstyle{\thelstnumber}}\kern\lst@numbersep}}\\%
     right:\def\lst@PlaceNumber{\rlap{\normalfont
                \kern\linewidth \kern\lst@numbersep
                \UseTaggingSocket{listings/number}{}{\lst@numberstyle{\thelstnumber}}}}%
    }{\PackageError{Listings}{value `#1' for keyword `numbers' not supported}\@ehc}}
\lst@AddToHook{PreInit}{\ifhmode \UseTaggingSocket{inline/begin}{tag=\UseStructureName{text/verb}}\fi}
\lst@AddToHook{DeInit}{\ifhmode \UseTaggingSocket{inline/end}%
    \else \edef\@current@float@struct{\UseName{tag_get:n}{struct_num}}\BlockEnvEnd \fi}
\def\lst@makecaption{\AssignStructureRole{para/textblock}{Span}\def\@captype{lstlisting}\@makecaption}
\def\lst@maketitle{\def\@captype{lstlisting}%
    \AssignStructureRole{para/textblock}{Span}%
    \AssignTaggingSocketPlug{caption/label/begin}{noop}%
    \AssignTaggingSocketPlug{caption/label/end}{noop}%
    \AssignSocketPlug{caption/label}{noop}%
    \@makecaption\lst@title@dropdelim}
\fi
%% The spacing, as without tagging (both releases): listings adds its own skips and break points
%% (\lstset aboveskip, belowskip), so the block around a listing adds none ...
\DeclareInstance{block}{lstlisting-1}{std}{%
    left-margin = 0pt, para-vspace = 0pt, begin-vspace = 0pt, end-vspace = 0pt,
    begin-extra-vspace = 0pt, end-extra-vspace = 0pt, begin-unchained-vspace = 0pt,
    begin-penalty = 10000, end-penalty = 10000}
\DeclareInstanceCopy{block}{lstlisting-2}{lstlisting-1}
\DeclareInstanceCopy{block}{lstlisting-3}{lstlisting-1}
\DeclareInstanceCopy{block}{lstlisting-4}{lstlisting-1}
\DeclareInstanceCopy{block}{lstlisting-5}{lstlisting-1}
\DeclareInstanceCopy{block}{lstlisting-6}{lstlisting-1}
%% ... no stretch above it (the block adds the paragraph skip), and no skips of nothing (they stayed at
%% the foot of a page that breaks at the listing)
\def\lst@EveryDisplay{%
    \MathCollectFalse
    \parskip\z@skip
    \ifx\lst@float\relax \SimpleBlockEnv{lstlisting}{}%
    \else \SimpleBlockEnv{lstlisting}{tagging-recipe = standalone}\fi
    \thesis@unskipzero}
\def\thesis@unskipzero{\let\thesis@next\relax
    \ifvmode\ifnum\lastnodetype=11 % (a skip)
        \thesis@ifzeroskip\lastskip{\unskip\let\thesis@next\thesis@unskipzero}{}\fi\fi
    \thesis@next}

%% Spaces in code as space characters (both releases): listings sets each word in its own box, so the
%% space between two words is not next to a letter and LaTeX's tagging adds no space character there.
%% The tagged text read "importnumpyasnp"; with \pdffakespace (a space character, no width) it reads
%% "import numpy as np". The pages do not change. (Indentation stays without characters.)
\def\lst@outputspace{\pdffakespace\ }

\ExplSyntaxOn
%% ---------------------------------------------------------------------------------------------
%% 2. Equation numbers in the MathML (only when the MathML is made: luamml, math/setup in
%% \DocumentMetadata, with the OpenType maths fonts).
%%   multline: luamml leaves multline out of the environments whose equation number gets its own
%%   structure, but still refers to that structure when it writes the MathML ("structure with label 0
%%   is unknown"): treat multline like equation and gather.
%%   equation with split: LaTeX keeps only the split's table in the MathML and drops the equation
%%   number; it is kept back here and added again around that table (tagging.lua).
%%   tagging.lua also makes a number one <mn> (siunitx's 6.022\,140\,76 was five pieces), a word in
%%   upright letters one <mi> (\mathrm{disc}), and leaves out boxes with nothing to read.
\AddToHook{begindocument}
  {
    \IfPackageLoadedT{luamml}
      {
        \tl_if_exist:cF { c__luamml_label_multline_tl }
          { \tl_const:cn { c__luamml_label_multline_tl } { } }
        \directlua{dofile(kpse.find_file('template/tagging.lua'))}
        \NewTaggingSocketPlug{math/luamml/mtable/tag/save}{thesis}
          {
            \bool_if:NTF \l__math_gathersplit_bool
              { \use:c{thesis@luamml@savetag} \l__luamml_amsmath_tag_struct_num_int }
              { \__luamml_amsmath_save_tag_with_struct_elem:N \l__luamml_amsmath_tag_struct_num_int }
          }
        \AssignTaggingSocketPlug{math/luamml/mtable/tag/save}{thesis}
        \NewTaggingSocketPlug{math/luamml/mtable/tag/set}{thesis}
          { \bool_if:NF \l__math_gathersplit_bool { \__luamml_amsmath_set_tag: } }
        \AssignTaggingSocketPlug{math/luamml/mtable/tag/set}{thesis}
        \NewTaggingSocketPlug{math/luamml/mtable/innertable/finalize}{thesis}
          {
            \__luamml_amsmath_finalize_inner_table:
            \bool_if:NT \l__math_gathersplit_bool { \use:c{thesis@luamml@wraptag} }
          }
        \AssignTaggingSocketPlug{math/luamml/mtable/innertable/finalize}{thesis}
      }
  }

%% ---------------------------------------------------------------------------------------------
%% 3. Theorems, proofs and other blocks that stand on their own (LaTeX's code of June 2026, which has
%% the socket block/endpe; the November 2026 code needs nothing here). That code ends them without the
%% usual note that what follows straight on continues the paragraph (\if@endpe), which its tagging
%% needs, but it changes the look: text straight after was indented, and a theorem, proof or list
%% straight after got \partopsep more space above it. Here the look stays as in the untagged editions,
%% while the tagging code still sees what it expects.
\socket_if_exist:nT { block/endpe }
  {
    \bool_new:N \g__thesis_block_endpe_bool   % a block on its own has just ended
    \bool_new:N \l__thesis_block_pre_bool
    \bool_new:N \l__thesis_block_space_bool
    \bool_new:N \l__thesis_block_after_bool
    \cs_new_protected:Npn \__thesis_block_after_standalone:
      {
        \bool_gset_true:N \g__thesis_block_endpe_bool
        \def \par
          {
            \use:c { @restorepar } \bool_gset_false:N \g__thesis_block_endpe_bool
            \everypar { } \par
          }
        \everypar
          {
            { \setbox 0 \lastbox }
            \everypar { } \bool_gset_false:N \g__thesis_block_endpe_bool
          }
      }
    \socket_new_plug:nnn { block/endpe } { thesis-standalone }
      { \legacy_if_gset_false:n { @endpe } \aftergroup \__thesis_block_after_standalone: }
    \AddToHook{begindocument}
      {
        \cs_new_eq:NN \__thesis_block_recipe_standalone: \__block_recipe_standalone:
        \cs_gset_protected:Npn \__block_recipe_standalone:
          {
            \__thesis_block_recipe_standalone:
            \socket_assign_plug:nn { block/endpe } { thesis-standalone }
          }
        \NewTaggingSocketPlug { block/recipe } { thesis }
          {
            \legacy_if:nTF { @endpe }
              { \bool_set_true:N \l__thesis_block_pre_bool }
              { \bool_set_false:N \l__thesis_block_pre_bool }
            \bool_set_eq:NN \l__thesis_block_space_bool \l__thesis_block_pre_bool
            \bool_if:NT \g__thesis_block_endpe_bool
              { \bool_set_true:N \l__thesis_block_space_bool }
            \bool_gset_false:N \g__thesis_block_endpe_bool
            \use:c { __block_recipe_#1: }
            \legacy_if:nTF { @endpe }
              { \bool_set_true:N \l__thesis_block_after_bool }
              { \bool_set_false:N \l__thesis_block_after_bool }
            \bool_if:NT \l__thesis_block_space_bool
              {
                % the space above as without tagging, then the state the tagging code set
                \legacy_if_gset_true:n { @endpe }
                \cs_set_eq:NN \__thesis_displayblock_begin: \__kernel_displayblock_begin:
                \cs_set_protected:Npn \__kernel_displayblock_begin:
                  {
                    \bool_if:NTF \l__thesis_block_after_bool
                      { \legacy_if_gset_true:n { @endpe } }
                      { \legacy_if_gset_false:n { @endpe } }
                    \__thesis_displayblock_begin:
                  }
              }
          }
        \AssignTaggingSocketPlug { block/recipe } { thesis }
      }
  }

%% ---------------------------------------------------------------------------------------------
%% 4. Figures side by side (a minipage each, with its own caption, in one figure): a figure may have only
%% one caption in a tagged PDF, so each minipage after the first starts a figure of its own in the
%% structure (the pages are unchanged).
\int_new:N \g__thesis_float_captions_int
\AddToHook{cmd/@floatboxreset/after}{ \int_gzero:N \g__thesis_float_captions_int }
\AddToHook{cmd/end@float/before}{ \int_gzero:N \g__thesis_float_captions_int }
\AddToHook{cmd/end@dblfloat/before}{ \int_gzero:N \g__thesis_float_captions_int }
\AddToHook{cmd/caption/before}{ \int_gincr:N \g__thesis_float_captions_int }
\AddToHook{env/minipage/before}
  {
    \bool_lazy_and:nnT
      { \cs_if_exist_p:c { @captype } }
      { \int_compare_p:nNn { \g__thesis_float_captions_int } > { 0 } }
      {
        \tagpdfsetup{float/split}
        \int_gzero:N \g__thesis_float_captions_int
      }
  }
\ExplSyntaxOff

%% ---------------------------------------------------------------------------------------------
%% 5. tagpdf stores the number of pages, marked-content items and structure elements of the last compile
%% (the label @tag@LastPage). The two totals change whenever anything earlier in the thesis gains or
%% loses a tagged item, such as the table of contents once the references are in it, and LaTeX would then
%% warn that labels may have changed and ask for another compile, although tagpdf (in LuaLaTeX) does not
%% use them for the PDF. Only a change in the number of pages counts here, as for any other label.
%% (latexmkrc ignores the same line.)
\ExplSyntaxOn
\cs_new:Npn \__thesis_abspage:w #1#2#3 \q_stop { \tl_if_eq:nnT {#1} {abspage} {#2} }
\AddToHook{enddocument/afterlastpage}
  {
    \cs_set_eq:NN \__thesis_label_testdef:nn \@kernel@new@label@record@testdef
    \cs_set_protected:Npn \@kernel@new@label@record@testdef #1#2
      {
        \str_if_eq:nnTF {#1} {@tag@LastPage}
          {
            \tl_if_exist:cTF { r@#1 }
              {
                \str_if_eq:eeF
                  { \exp_last_unbraced:Nv \__thesis_abspage:w { r@#1 } {} {} {} \q_stop }
                  { \__thesis_abspage:w #2 {} {} {} \q_stop }
                  { \@tempswatrue }
              }
              { \@tempswatrue }
          }
          { \__thesis_label_testdef:nn {#1} {#2} }
      }
  }
\ExplSyntaxOff

%% ---------------------------------------------------------------------------------------------
%% 6. Table cells: LaTeX gives each header cell its scope (column, row or both), and each cell that
%% spans columns its span, through attribute classes (/C /TH-col, /C /colspan-2, defined once in the
%% ClassMap). That is valid in PDF 1.7 and 2.0, and veraPDF reads it, but PAC does not: it reports
%% every header cell as having "no associated subcells", and every table with a \multicolumn as not
%% regular (tested on plain LaTeX documents in PDF 1.7 and 2.0). So the table attributes of a cell's
%% classes are also written on the cell itself (/A << /O /Table ... >>), next to the classes, when the
%% structure is written out at the end. Remove once PAC reads attribute classes, or once LaTeX writes
%% them on the cell itself (the package watch fingerprints \__tag_struct_write_obj:n, which this wraps).
\ExplSyntaxOn
\cs_if_exist:NT \__tag_struct_write_obj:n
  {
    \cs_new_eq:NN \__thesis_struct_write_obj:n \__tag_struct_write_obj:n
    \cs_set_protected:Npn \__tag_struct_write_obj:n #1
      {
        \prop_get:cnNT { g__tag_struct_#1_prop } { C } \l__thesis_tmpa_tl
          {
            % the classes: one name, or several in [ ]
            \str_set:Ne \l__thesis_tmpa_str { \tl_to_str:N \l__thesis_tmpa_tl }
            \str_remove_all:Nn \l__thesis_tmpa_str { [ }
            \str_remove_all:Nn \l__thesis_tmpa_str { ] }
            \seq_set_split:NnV \l__thesis_tmpa_seq { ~ } \l__thesis_tmpa_str
            % the entries of those owned by Table, without their /O /Table
            \str_clear:N \l__thesis_tmpb_str
            \seq_map_inline:Nn \l__thesis_tmpa_seq
              {
                \prop_get:NnNT \g__tag_attr_entries_prop {##1} \l__thesis_tmpb_tl
                  {
                    \str_set:Ne \l__thesis_tmpc_str { \tl_to_str:N \l__thesis_tmpb_tl }
                    \str_if_in:NnT \l__thesis_tmpc_str { /Table }
                      {
                        \seq_set_split:NnV \l__thesis_tmpb_seq { ~ } \l__thesis_tmpc_str
                        \seq_map_inline:Nn \l__thesis_tmpb_seq
                          {
                            \str_if_eq:nnF {####1} { /O }
                              {
                                \str_if_eq:nnF {####1} { /Table }
                                  { \tl_if_blank:nF {####1} { \str_put_right:Nn \l__thesis_tmpb_str { ~ ####1 } } }
                              }
                          }
                      }
                  }
              }
            \str_if_empty:NF \l__thesis_tmpb_str
              {
                \prop_get:cnNTF { g__tag_struct_#1_prop } { A } \l__thesis_tmpa_tl
                  { % (other attributes already: one array of them all, without nesting arrays)
                    \tl_set:Ne \l__thesis_tmpa_tl { \tl_trim_spaces:V \l__thesis_tmpa_tl }
                    \str_if_eq:eeT { \tl_head:N \l__thesis_tmpa_tl } { [ }
                      { \tl_set:Ne \l__thesis_tmpa_tl { \tl_range:Nnn \l__thesis_tmpa_tl { 2 } { -2 } } }
                    \prop_gput:cne { g__tag_struct_#1_prop } { A }
                      { [ \l__thesis_tmpa_tl \c_space_tl << /O /Table \l__thesis_tmpb_str \c_space_tl >> ] }
                  }
                  {
                    \prop_gput:cne { g__tag_struct_#1_prop } { A }
                      { << /O /Table \l__thesis_tmpb_str \c_space_tl >> }
                  }
              }
          }
        \__thesis_struct_write_obj:n {#1}
      }
  }
\tl_new:N \l__thesis_tmpa_tl
\tl_new:N \l__thesis_tmpb_tl
\str_new:N \l__thesis_tmpa_str
\str_new:N \l__thesis_tmpb_str
\str_new:N \l__thesis_tmpc_str
\seq_new:N \l__thesis_tmpa_seq
\seq_new:N \l__thesis_tmpb_seq
\ExplSyntaxOff

\makeatother
\fi
        % tagged PDF: what LaTeX's tagging code does not cover yet (listings; before
                                %  their style, which it would reset)
%% Code listings (listings package).
%% Digital edition: coloured syntax highlighting. Print editions: black and white (bold keywords,
%% grey italic comments), because printers often charge every page with colour as a colour page.
%% Colours with enough contrast against white for small text (4.5:1 or more): green 5.6, grey 5.7,
%% purple 6.6, blue 8.6.
\definecolor{codegreen}{rgb}{0,0.478,0}
\definecolor{codegrey}{gray}{0.4}
\definecolor{codepurple}{rgb}{0.58,0,0.82}
\lstset{
    frame=tb,
    language=Python,
    aboveskip=3mm,
    belowskip=3mm,
    showstringspaces=false,
    columns=flexible,
    basicstyle={\small\ttfamily},
    numbers=left,
    numbersep=5pt,
    numberstyle=\tiny\color{codegrey},
    breaklines=true,
    breakatwhitespace=false,
    tabsize=3,
    extendedchars=true,
    literate={á}{{\'a}}1 {ã}{{\~a}}1 {é}{{\'e}}1 {ü}{{\"u}}1,
}
\ifprintversion
    \lstset{
        keywordstyle=\bfseries,
        commentstyle=\itshape\color{codegrey},
        stringstyle=\color{black},
    }
\else
    \lstset{
        keywordstyle=\color{blue},
        commentstyle=\color{codegreen},
        stringstyle=\color{codepurple},
    }
\fi
     % their style: colour on screen, black and white in print

%% Todo notes while writing: \todo{...} in the margin, \todo[inline]{...} in the text,
%% \missingfigure{...} as a placeholder figure. All of them disappear in the final version (version=final).
\usepackage[textsize=scriptsize, color=tudelft-yellow, linecolor=black!50]{todonotes}
%% Changes for a supervisor to review: \added{new text}, \deleted{old text}, \replaced{new}{old},
%% \highlight{text}, \comment{remark}; \listofchanges lists them. In the final version (version=final) the
%% marks disappear and only the new text remains.
\usepackage{changes}            % https://ctan.org/pkg/changes

%% Paragraphs: book style, first line indented and no blank space between paragraphs.
%% (The first paragraph after a heading is not indented, as usual.)
\setlength{\parindent}{1.5em}
\setlength{\parskip}{0pt plus 1pt} % a tiny stretch, so a page that breaks early spreads the gap over all paragraphs

%% Hyperlinks and PDF bookmarks
\usepackage{bookmark}           % https://ctan.org/pkg/bookmark %% PDF bookmarks (loads hyperref)
%% Hyperlinks: invisible in the print editions (they would only show up as coloured text on paper),
%% brick red in the digital edition (the colour "link", set in settings/colours.tex).
\ifprintversion
    \hypersetup{hidelinks}
\else
    \hypersetup{colorlinks=true, linkcolor=link, citecolor=link, urlcolor=link}
\fi
%% Bookmarks: numbered ("2.1 Text"), and one level deeper than the contents (down to subsections),
%% so the PDF can be navigated in more detail than the printed contents show.
\bookmarksetup{numbered, depth=2}
%% PDF metadata (title, author, keywords) comes from settings/thesisinfo.tex (template/pdfa.tex).

%% Nomenclature (acronyms and symbols), after hyperref (its entries link to their rows)
%% =====================================================================
%% Nomenclature: acronyms and symbols, defined once and listed automatically.
%% Entries are defined in settings/nomenclature.tex (preamble) and printed in chapters/00Nomenclature.tex.
%% In the text:  \gls{bem}    -> "blade element momentum theory (BEM)" the first time, "BEM" after that
%%               \Gls{bem}    -> the same with a capital at the start of a sentence
%%               \glspl{bem}  -> plurals: "... theories (BEMs)", then "BEMs"; \Glspl likewise
%%               \gls{rho}    -> the symbol, in text or maths
%%               \glsfmtshort{bem}, \glsfmtlong{bem} -> always the short or long form (headings, captions)
%%               \glsreset{bem}, \glsresetall -> write out in full again at the next use
%% Every entry links to its row in the nomenclature. Only entries used in the text are listed,
%% unless \glsaddall is given (chapters/00Nomenclature.tex); used entries are noted in the .aux
%% file, so a new one appears after the next run.
%% (The template's own code: the glossaries package did the same but cost about 6 s per run.)
%% =====================================================================
\makeatletter
\ExplSyntaxOn
%% Every entry keeps its fields in token lists \g__thesis_nom_<label>_<field>_tl; the labels of
%% each list are kept in definition order in \g__thesis_nom_<type>_seq.
\cs_new:Npn \__thesis_nom:nn #1#2 { g__thesis_nom_ #1 _ #2 _tl }
\cs_new_protected:Npn \__thesis_nom_set:nnn #1#2#3 { \tl_gclear_new:c { \__thesis_nom:nn {#1} {#2} } \tl_gset:cn { \__thesis_nom:nn {#1} {#2} } {#3} }
\cs_new:Npn \__thesis_nom_get:nn #1#2 { \tl_if_exist:cT { \__thesis_nom:nn {#1} {#2} } { \tl_use:c { \__thesis_nom:nn {#1} {#2} } } }
%% (the same, not expanded further in an e-expansion: fields hold maths such as \vec{u})
\cs_new:Npn \__thesis_nom_raw:nn #1#2 { \tl_if_exist:cT { \__thesis_nom:nn {#1} {#2} } { \exp_not:v { \__thesis_nom:nn {#1} {#2} } } }
\clist_map_inline:nn { acronyms, dimless, latin, greek, super, sub, indices, notation } { \seq_new:c { g__thesis_nom_#1_seq } }
\bool_new:N \g__thesis_nom_all_bool

%% #1 type, #2 sort key (expanded), #3 label, then the fields as a key = value list
\cs_new_protected:Npn \__thesis_nom_new:nnnn #1#2#3#4
  {
    \tl_if_exist:cT { \__thesis_nom:nn {#3} { type } }
      { \msg_error:nnn { thesis } { nom-duplicate } {#3} }
    \seq_gput_right:cn { g__thesis_nom_#1_seq } {#3}
    %% sub- and superscripts also go, in the order they are defined, into one list of indices
    \str_case:nn {#1} { {sub} { \seq_gput_right:Nn \g__thesis_nom_indices_seq {#3} } {super} { \seq_gput_right:Nn \g__thesis_nom_indices_seq {#3} } }
    \__thesis_nom_set:nnn {#3} { type } {#1}
    \__thesis_nom_set:nnn {#3} { sort } {#2}
    \keyval_parse:nnn { \use_none:n } { \__thesis_nom_set:nnn {#3} } {#4}
    \bool_new:c { g__thesis_nom_#3_used_bool }
    \bool_new:c { g__thesis_nom_#3_listed_bool }
  }
\msg_new:nnn { thesis } { nom-duplicate } { Nomenclature~entry~'#1'~defined~twice. }
\msg_new:nnn { thesis } { nom-undefined } { Nomenclature~entry~'#1'~not~defined~(in~settings/nomenclature.tex). }

%% ---------------------------------------------------------------------
%% Commands for defining entries (use them in settings/nomenclature.tex).
%% The optional first argument is a sort key. By default, acronyms are sorted by their short form and
%% Latin symbols and dimensionless numbers by their label (case ignored, lower case first); Greek symbols
%% and sub- and superscripts keep the order in which you define them.
%% ---------------------------------------------------------------------
%% \nomacronym[sort]{label}{SHORT}{long form}   (abbreviations of any kind: BEM, RPM, ...)
\NewDocumentCommand \nomacronym { O{} m m m }
  { \thesis@withsort {#1} {#3} { \__thesis_nom_new:nnnn { acronyms } } {#2} { short = {#3}, long = {#4} } }
%% \nomdimless[sort]{label}{symbol}{description}{definition}       (symbol and definition: maths, U/a)
\NewDocumentCommand \nomdimless { O{} m m m m }
  { \thesis@withsort {#1} {#2} { \__thesis_nom_new:nnnn { dimless } } {#2} { name = {#3}, desc = {#4}, definition = {#5} } }
%% \nomlatin[sort]{label}{symbol}{description}{dimensions}{units}   (symbol and dimensions: maths)
\NewDocumentCommand \nomlatin { O{} m m m m m }
  { \thesis@withsort {#1} {#2} { \__thesis_nom_new:nnnn { latin } } {#2} { name = {#3}, desc = {#4}, dims = {#5}, units = {#6} } }
%% \nomgreek[sort]{label}{symbol}{description}{dimensions}{units}
\NewDocumentCommand \nomgreek { O{} m m m m m }
  { \thesis@withsort {#1} {#2} { \__thesis_nom_new:nnnn { greek } } {#2} { name = {#3}, desc = {#4}, dims = {#5}, units = {#6} } }
%% \nomsuper[sort]{label}{symbol}{description}
\NewDocumentCommand \nomsuper { O{} m m m }
  { \thesis@withsort {#1} {#2} { \__thesis_nom_new:nnnn { super } } {#2} { name = {#3}, desc = {#4} } }
%% \nomsub[sort]{label}{symbol}{description}
\NewDocumentCommand \nomsub { O{} m m m }
  { \thesis@withsort {#1} {#2} { \__thesis_nom_new:nnnn { sub } } {#2} { name = {#3}, desc = {#4} } }
%% \nomnotation[sort]{label}{symbol}{description}   (operators and conventions: \vec{u}, \grad, (\cdot)\transpose)
\NewDocumentCommand \nomnotation { O{} m m m }
  { \thesis@withsort {#1} {#2} { \__thesis_nom_new:nnnn { notation } } {#2} { name = {#3}, desc = {#4} } }

%% \thesis@withsort{sort key}{label}{\cmd}: calls \cmd{key}, the key made (expanded) from the sort key
%% given, or from the label when that is empty.
\cs_new_protected:Npn \thesis@withsort #1#2#3
  { \tl_if_blank:nTF {#1} { \exp_args:Ne #3 { \thesis@sortkey {#2} } } { \exp_args:Ne #3 { \thesis@sortkey {#1} } } }
%% Sort keys ignore case, and a lower-case letter comes before the same capital (f before F, t before T),
%% the order of a dictionary: "f" -> "f!0" and "F" -> "f!1".
\cs_new:Npn \thesis@sortkey #1
  { \str_lowercase:n {#1} ! \str_map_function:nN {#1} \__thesis_caseflag:N }
\cs_new:Npn \__thesis_caseflag:N #1
  { \str_if_eq:eeTF {#1} { \str_lowercase:n {#1} } { 0 } { 1 } }

%% ---------------------------------------------------------------------
%% Use in the text
%% ---------------------------------------------------------------------
%% Symbols and acronyms in the text stay in the text colour, also in the digital edition: they still
%% link to their nomenclature entry, but colour is kept for references (sections, equations, figures,
%% sources, URLs), so maths is never coloured.
\cs_new_protected:Npn \__thesis_nom_link:nn #1#2 % label, text
  { \hyperlink { glo:#1 } { \normalcolor #2 } }
\cs_new_protected:Npn \__thesis_nom_check:nT #1#2
  {
    \tl_if_exist:cTF { \__thesis_nom:nn {#1} { type } }
      {
        \bool_if:cF { g__thesis_nom_#1_listed_bool }
          {
            \bool_gset_true:c { g__thesis_nom_#1_listed_bool }
            \iow_now:ce { @auxout } { \token_to_str:N \thesis@nomused { #1 } }
          }
        #2
      }
      { \msg_error:nnn { thesis } { nom-undefined } {#1} \textbf { ?? } }
  }
%% #1 label, #2 capital (true/false), #3 plural (true/false)
\cs_new_protected:Npn \__thesis_gls:nnn #1#2#3
  {
    \__thesis_nom_check:nT {#1}
      {
        \str_if_eq:eeTF { \__thesis_nom_get:nn {#1} { type } } { acronyms }
          {
            \bool_if:cTF { g__thesis_nom_#1_used_bool }
              {
                \thesisexpansion { \__thesis_nom_long:nnn {#1} { \c_false_bool } {#3} } % (for screen readers)
                  { \__thesis_nom_link:nn {#1} { \__thesis_nom_short:nn {#1} {#3} } }
              }
              {
                \bool_gset_true:c { g__thesis_nom_#1_used_bool }
                \__thesis_nom_link:nn {#1}
                  { \__thesis_nom_long:nnn {#1} {#2} {#3} ~ ( \__thesis_nom_short:nn {#1} {#3} ) }
              }
          }
          { \__thesis_nom_link:nn {#1} { \ensuremath { \__thesis_nom_get:nn {#1} { name } } } }
      }
  }
\cs_new:Npn \__thesis_nom_short:nn #1#2 { \__thesis_nom_get:nn {#1} { short } \bool_if:nT {#2} { s } }
\cs_new:Npn \__thesis_nom_long:nnn #1#2#3
  {
    \bool_if:nTF {#2}
      { \text_titlecase_first:e { \__thesis_nom_get:nn {#1} { long } } }
      { \__thesis_nom_get:nn {#1} { long } }
    \bool_if:nT {#3} { s }
  }
\cs_generate_variant:Nn \text_titlecase_first:n { e }
\NewDocumentCommand \gls { m } { \__thesis_gls:nnn {#1} { \c_false_bool } { \c_false_bool } }
\NewDocumentCommand \Gls { m } { \__thesis_gls:nnn {#1} { \c_true_bool } { \c_false_bool } }
\NewDocumentCommand \glspl { m } { \__thesis_gls:nnn {#1} { \c_false_bool } { \c_true_bool } }
\NewDocumentCommand \Glspl { m } { \__thesis_gls:nnn {#1} { \c_true_bool } { \c_true_bool } }
%% always the short or the long form, not counted as a use (headings, captions, the list of figures)
\NewDocumentCommand \glsfmtshort { m } { \__thesis_nom_check:nT {#1} { \__thesis_nom_short:nn {#1} { \c_false_bool } } }
\NewDocumentCommand \glsfmtlong { m } { \__thesis_nom_check:nT {#1} { \__thesis_nom_long:nnn {#1} { \c_false_bool } { \c_false_bool } } }
\NewDocumentCommand \glsreset { m } { \bool_if_exist:cT { g__thesis_nom_#1_used_bool } { \bool_gset_false:c { g__thesis_nom_#1_used_bool } } }
\NewDocumentCommand \glsresetall { }
  {
    \clist_map_inline:nn { acronyms, dimless, latin, greek, super, sub, notation }
      { \seq_map_inline:cn { g__thesis_nom_##1_seq } { \bool_gset_false:c { g__thesis_nom_####1_used_bool } } }
  }
%% list every entry, used or not
\NewDocumentCommand \glsaddall { } { \bool_gset_true:N \g__thesis_nom_all_bool }
%% (in the .aux file: the entries used in the text in the previous run)
\cs_new_protected:Npn \thesis@nomused #1 { \tl_gset:cn { g__thesis_nom_#1_inaux_tl } { } }
\ExplSyntaxOff

%% Dimensions are given in maths ("\Lunit\Tunit^{-1}"); "--" gives an en dash for dimensionless quantities.
\newcommand{\thesis@dims}[1]{%
    \def\thesis@arg{#1}\def\thesis@dash{--}%
    \ifx\thesis@arg\thesis@dash \textendash \else \ensuremath{#1}\fi}

%% MLT dimensions (https://www.ldeo.columbia.edu/~martins/hydro/case_studies/units_dimensions.html)
\newcommand{\Lunit}{\mathrm{L}}   % length
\newcommand{\Munit}{\mathrm{M}}   % mass
\newcommand{\Tunit}{\mathrm{T}}   % time
\newcommand{\Qunit}{\Theta}       % temperature (ISQ dimension symbol)

%% ---------------------------------------------------------------------
%% Layout of the lists
%% ---------------------------------------------------------------------
%% A hairline (0.3 pt, like the rule under the running heads) separates the column headers.
%% Narrow symbol, dimensions and units columns and 4 pt column gaps leave the description as much
%% room as possible, also on the narrow demy page.
\newlength\symbolwidth      \setlength\symbolwidth{4em}
\newlength\dimensionswidth  \setlength\dimensionswidth{5em}
\newlength\unitswidth       \setlength\unitswidth{5.5em}
\newlength\descriptionwidth

\newcommand{\nomheader}[1]{{\sffamily\mbox{#1}}} % Source Sans Regular; \mbox: no hyphenation in narrow columns
%% List titles look like section headings: semibold, 1.15 x text size, the same space around them;
%% a title only starts a list where its header and first rows fit below it.
%% (\needspace's stretchable glue confuses longtable, which may then set a table's header above the title)
\newcommand{\thesis@needspace}[1]{\par\ifdim\dimexpr\pagegoal-\pagetotal\relax<#1\relax\newpage\fi}
\newcommand{\nomtitle}[1]{\par\addvspace{1.5\onelineskip}\thesis@needspace{7\onelineskip}\noindent{\thesis@secsize\bfseries\headingfont #1}\par\nopagebreak\vspace{0.5\onelineskip}}

%% The rows of a list are collected first and set as one longtable at the end, so that the first
%% two and the last two rows can be kept together: a page break never leaves a single row under the
%% header at the bottom of a page, or alone at the top of the next.
\ExplSyntaxOn
\seq_new:N \g__thesis_nom_seq
\tl_new:N \l__thesis_nom_tl
\cs_new_protected:Npn \thesis@nomrow #1 { \seq_gput_right:Nn \g__thesis_nom_seq {#1} }
\cs_new_protected:Npn \thesis@nomtable #1#2 % #1 column spec, #2 header row
  {
    \tl_clear:N \l__thesis_nom_tl
    \int_set:Nn \l_tmpa_int { \seq_count:N \g__thesis_nom_seq }
    \seq_map_indexed_inline:Nn \g__thesis_nom_seq
      {
        \tl_put_right:Nn \l__thesis_nom_tl {##2}
        \bool_lazy_or:nnTF
          { \int_compare_p:nNn {##1} = { 1 } }
          { \int_compare_p:nNn {##1} = { \l_tmpa_int - 1 } }
          { \tl_put_right:Nn \l__thesis_nom_tl { \tabularnewline* } }
          { \tl_put_right:Nn \l__thesis_nom_tl { \tabularnewline } }
      }
    \seq_gclear:N \g__thesis_nom_seq
    \exp_args:NnnV \__thesis_nom_longtable:nnn {#1} {#2} \l__thesis_nom_tl
  }
\cs_new_protected:Npn \__thesis_nom_longtable:nnn #1#2#3
  { \begin{longtable}{#1} #2 \tabularnewline* \midrule[0.3pt] \endhead #3 \end{longtable} }
\ExplSyntaxOff

%% ---------------------------------------------------------------------
%% The lists
%% ---------------------------------------------------------------------
\ExplSyntaxOn
\seq_new:N \l__thesis_nom_labels_seq
%% sort orders: nocase (by sort key) or def (order of definition)
\cs_new_protected:Npn \__thesis_nom_sorted:nn #1#2 % type, order
  {
    \seq_clear:N \l__thesis_nom_labels_seq
    \seq_map_inline:cn { g__thesis_nom_#1_seq }
      {
        \bool_lazy_or:nnT { \g__thesis_nom_all_bool } { \tl_if_exist_p:c { g__thesis_nom_##1_inaux_tl } }
          { \seq_put_right:Nn \l__thesis_nom_labels_seq {##1} }
      }
    \str_if_eq:nnT {#2} { nocase }
      {
        \seq_sort:Nn \l__thesis_nom_labels_seq
          {
            \str_compare:eNeTF { \__thesis_nom_get:nn {##1} { sort } } > { \__thesis_nom_get:nn {##2} { sort } }
              { \sort_return_swapped: } { \sort_return_same: }
          }
      }
  }
%% the cells after Symbol | Description for each style
\cs_new:Npn \__thesis_nom_cells_plain:n #1 { }
\cs_new:Npn \__thesis_nom_cells_definition:n #1
  { & \exp_not:N \ensuremath { \__thesis_nom_raw:nn {#1} { definition } } }
\cs_new:Npn \__thesis_nom_cells_units:n #1 { & \__thesis_nom_raw:nn {#1} { units } }
\cs_new:Npn \__thesis_nom_cells_dimensions:n #1
  {
    & \exp_not:N \thesis@dims { \__thesis_nom_raw:nn {#1} { dims } }
    & \__thesis_nom_raw:nn {#1} { units }
  }
%% #5 sets the table (\__thesis_nom_table:nnn with the columns after Description)
\cs_new_protected:Npn \__thesis_nom_print:nnnnn #1#2#3#4#5 % type, order, style, title, table
  {
    \__thesis_nom_sorted:nn {#1} {#2}
    \seq_if_empty:NF \l__thesis_nom_labels_seq
      {
        \nomtitle {#4}
        \group_begin:
          \setlength\tabcolsep{4pt}
          \setlength\LTleft{0pt}\setlength\LTright{0pt}\setlength\LTpre{0pt}\setlength\LTpost{0pt}
          \seq_map_inline:Nn \l__thesis_nom_labels_seq
            {
              \str_if_eq:nnTF {#3} { acronyms }
                {
                  \exp_args:Ne \thesis@nomrow
                    {
                      \exp_not:N \__thesis_nom_target:n {##1} { \__thesis_nom_raw:nn {##1} { short } }
                      & \exp_not:N \text_titlecase_first:n { \__thesis_nom_raw:nn {##1} { long } }
                    }
                }
                {
                  \exp_args:Ne \thesis@nomrow
                    {
                      \exp_not:N \__thesis_nom_target:n {##1}
                        { \exp_not:N \ensuremath { \__thesis_nom_symbol:nn {#1} {##1} } }
                      & \__thesis_nom_raw:nn {##1} { desc }
                      \use:c { __thesis_nom_cells_#3:n } {##1}
                    }
                }
            }
          #5
        \group_end:
      }
  }
%% the symbol in a list; in the list of indices with a dot for what it is attached to: (·)_∞, (·)^n
\cs_new:Npn \__thesis_nom_symbol:nn #1#2 % list, label
  {
    \str_if_eq:nnTF {#1} { indices }
      {
        { ( \exp_not:N \cdot ) }
        \str_if_eq:eeTF { \__thesis_nom_get:nn {#2} { type } } { super }
          { \exp_not:N \c_math_superscript_token } { \exp_not:N \c_math_subscript_token }
        { \__thesis_nom_raw:nn {#2} { name } }
      }
      { \__thesis_nom_raw:nn {#2} { name } }
  }
\cs_new_protected:Npn \__thesis_nom_target:n #1 { \leavevmode \hypertarget { glo:#1 } }
\cs_new_protected:Npn \__thesis_nom_table:nnn #1#2#3 % width of the extra columns, their spec, their headers
  {
    \setlength\descriptionwidth{\dimexpr\textwidth-\symbolwidth-2\tabcolsep-(#1)\relax}
    \thesis@nomtable{@{}p{\symbolwidth}>{\raggedright\arraybackslash}p{\descriptionwidth}#2@{}}
      {\nomheader{Symbol} & \nomheader{Description}#3}
  }
%% \printnom{list}{title}{style}{sort}: prints one list with its title (nothing if it is empty).
%% Lists: acronyms, dimless, latin, greek, notation, sub, super, and indices (sub- and superscripts in one).
%% Styles: nom-acronyms (Symbol | Description = long form), nom-plain (Symbol | Description),
%% nom-definition (... | Definition), nom-units (... | Units), nom-dimensions (... | Dimensions | Units).
\NewDocumentCommand \printnom { m m m m }
  {
    \str_case:nn {#3}
      {
        { nom-acronyms } { \__thesis_nom_print:nnnnn {#1} {#4} { acronyms } {#2} { \__thesis_nom_table:nnn { 0pt } { } { } } }
        { nom-plain } { \__thesis_nom_print:nnnnn {#1} {#4} { plain } {#2} { \__thesis_nom_table:nnn { 0pt } { } { } } }
        { nom-definition }
          {
            \__thesis_nom_print:nnnnn {#1} {#4} { definition } {#2}
              {
                \__thesis_nom_table:nnn { \dimensionswidth+\unitswidth+4\tabcolsep }
                  { p{\dimexpr\dimensionswidth+\unitswidth+2\tabcolsep} } { & \nomheader{Definition} }
              }
          }
        { nom-units }
          {
            \__thesis_nom_print:nnnnn {#1} {#4} { units } {#2}
              { \__thesis_nom_table:nnn { \unitswidth+2\tabcolsep } { p{\unitswidth} } { & \nomheader{Units} } }
          }
        { nom-dimensions }
          {
            \__thesis_nom_print:nnnnn {#1} {#4} { dimensions } {#2}
              {
                \__thesis_nom_table:nnn { \dimensionswidth+\unitswidth+4\tabcolsep }
                  { p{\dimensionswidth}p{\unitswidth} } { & \nomheader{Dimensions} & \nomheader{Units} }
              }
          }
      }
  }
\ExplSyntaxOff
\makeatother


%% Numbered statements, all in upright text with a bold label (amsthm, which must come before cleveref):
%%   formulation  "Formulation 2.1"  weak forms, optimisation problems, discrete schemes   \cref{form:...}
%%   definition, lemma, theorem, corollary  share one counter: "Definition 2.1", "Lemma 2.2", ...
%%   remark       "Remark."          unnumbered side note, italic label
%%   proof        "Proof."           ends with a square (\qedhere puts it on the last line of an equation)
%% Theorems are often set in italics; upright reads better when the statement is full of maths.
%% To rename one, change the word here and in its \crefname below.
%% =====================================================================
%% Numbered statements and proofs: amsthm (LaTeX's own code with \DocumentMetadata) and their look.
%% The statements themselves (\newtheorem) and their names are in template/packages.tex.
%% Read by template/packages.tex.
%% =====================================================================
\let\openbox\relax  % also defined by newtxmath; let amsthm define it
\usepackage{amsthm}
%% (With \DocumentMetadata, LaTeX's own code stands in for amsthm and does not define \openbox, the square
%% at the end of a proof: amsthm's definition.)
\providecommand{\openbox}{\leavevmode\hbox to.77778em{\hfil\vrule\vbox to.675em{\hrule width.6em\vfil\hrule}\vrule\hfil}}
%% Proofs stretch like the other statements (amsthm gives them 6pt plus 6pt, so they took up most of the
%% slack on a page that breaks early, leaving holes around them).
%% (With \DocumentMetadata, which the template requires, LaTeX's own code replaces amsthm: proofs are its
%% "proof" blocks.)
\EditInstance{block}{proof-1}{begin-vspace = 6pt plus 2pt}
\EditInstance{block}{proof-2}{begin-vspace = 6pt plus 2pt}
\EditInstance{captionedtext}{proof}{after-hspace = \labelsep, caption-unbreakable = true}% (as amsthm: a rigid separation, the label in a box)
\makeatother

%% The theorem style: head "Theorem 3.1." in bold; with a note, "Theorem 3.1 (Name)." with the note in regular weight.
%% (LaTeX's tagging code reads the last argument as a list of \thmname, \thmnumber and \thmnote only.)
\newtheoremstyle{thesis}{\topsep}{\topsep}{\normalfont}{0pt}{\bfseries}{.}{5pt plus 1pt minus 1pt}%
    {\thmname{#1}\thmnumber{ #2}\thmnote{ (#3)}}
%% (amsthm put only the space above and below a statement around it; LaTeX's own code adds \partopsep)
%% (and the paragraph skip of the first paragraph is part of the space above, as with amsthm)
\IfInstanceExistsT{block}{thm-thesis-1}{%
    \newskip\thesistheoremsep \thesistheoremsep=\topsep
    \newskip\thesisremarksep \thesisremarksep=\glueexpr\topsep/2\relax % (as amsthm: half the space around a theorem)
    \EditInstance{block}{thm-thesis-1}{begin-extra-vspace = 0pt, begin-vspace = \thesistheoremsep - \parskip, end-vspace = \thesistheoremsep}%
    \EditInstance{block}{thm-thesis-2}{begin-extra-vspace = 0pt, begin-vspace = \thesistheoremsep - \parskip, end-vspace = \thesistheoremsep}%
    \EditInstance{block}{thm-remark-1}{begin-vspace = \thesisremarksep - \parskip, end-vspace = \thesisremarksep}%
    \EditInstance{block}{thm-remark-2}{begin-vspace = \thesisremarksep - \parskip, end-vspace = \thesisremarksep}%
    \EditInstance{thmstyle}{remark}{title-format = #1}}% (as amsthm: no italic correction after "Remark", which LaTeX adds from November 2026)
   % amsthm, and the look of statements and proofs (template/statements.tex)
%% (the environments are defined after cleveref below, so it names shared-counter ones correctly)

%% Algorithms (pseudocode): "Algorithm 2.1", numbered per chapter, as a float with its caption on top.
%% Write the steps with \State, \If, \For, \While, \Repeat/\Until, \Return, \Comment: https://ctan.org/pkg/algorithmicx
%% List of algorithms: \listofalgorithms in main.tex (commented out there).
%% =====================================================================
%% Algorithms (pseudocode): the algorithm float and algpseudocode, with the look of listings and tables.
%% Read by template/packages.tex.
%% =====================================================================
\usepackage[chapter]{algorithm}
\usepackage{algpseudocode}
\makeatletter
%% A \label after \State refers to the line number. LaTeX's list code (November 2026) makes the label of an
%% item given in brackets the current label, so \Require, \Ensure and \Statex (\item[...]) would take
%% it over; algorithmicx sets the line number as the label only once. Set again after every item.
\algrenewcommand\ALG@beginalgorithmic{%
    \let\thesis@algitem\item
    \def\thesis@alglabel{\def\@currentlabel{\theALG@line}}%
    \def\thesis@algitem@opt[##1]{\thesis@algitem[##1]\thesis@alglabel}%
    \def\item{\@ifnextchar[\thesis@algitem@opt{\thesis@algitem\thesis@alglabel}}}
%% every algorithm numbers its lines from 1: give their link targets unique names
\renewcommand{\theHALG@line}{\thealgorithm.\arabic{ALG@line}}
%% Same look as code listings and tables: caption above (small, label in bold small capitals), then a
%% thin rule, the steps, and a thin rule.
\renewcommand\fs@ruled{\def\@fs@cfont{}\let\@fs@capt\floatc@ruled
    \def\@fs@pre{}\def\@fs@post{\kern2pt\hrule\relax}%
    \def\@fs@mid{\kern3pt\hrule\kern3pt}\let\@fs@iftopcapt\iftrue}
\floatstyle{ruled}\restylefloat{algorithm}
%% (with LaTeX's tagging code of November 2026, float.sty's caption ends the caption structure itself:
%% then it is begun here too)
\ExplSyntaxOn
\exp_args:Nee \str_if_in:nnTF { \token_to_meaning:N \float@caption } { \tl_to_str:n { caption/end } }
  { \def\thesis@floatcapbegin{\UseTaggingSocket{caption/begin}{\@current@float@struct}} }
  { \let\thesis@floatcapbegin\relax }
\ExplSyntaxOff
\renewcommand\floatc@ruled[2]{\thesis@floatcapbegin
    {\parindent\z@\thesis@capcentrefalse\tagpdfparaOff\leavevmode\thesis@captionbox{#1}{#2}\par\tagpdfparaOn}\@minipagefalse}
\makeatother
   % algorithm and algpseudocode, with the look of listings and tables
\algrenewcommand\algorithmicrequire{\textbf{Input:}}
\algrenewcommand\algorithmicensure{\textbf{Output:}}
\algrenewcommand\algorithmiccomment[1]{\hfill\textcolor{black!60}{\(\triangleright\) \textit{#1}}}

%% cleveref must be loaded last (after hyperref). Use \cref{fig:x} -> "Figure 2.1", \Cref at the start of a sentence.
\usepackage[capitalise,noabbrev,nameinlink]{cleveref}
\crefname{page}{page}{pages} % "on page 12", not "on Page 12"
\crefname{listing}{Listing}{Listings}
\crefname{formulation}{Formulation}{Formulations}
\crefname{definition}{Definition}{Definitions}
\crefname{lemma}{Lemma}{Lemmas}
\crefname{theorem}{Theorem}{Theorems}
\crefname{corollary}{Corollary}{Corollaries}
\crefname{algorithm}{Algorithm}{Algorithms}
\crefname{subfigure}{Figure}{Figures} % sub-figures and sub-tables (thesis.cls): "Figure 2.3a"
\crefname{subtable}{Table}{Tables}
%% Refer to an algorithm line with line~\ref{alg:...} (\cref would capitalise it mid-sentence).

%% The numbered statements described above (after cleveref, which needs that for shared counters)
\theoremstyle{thesis}
\newtheorem{formulation}{Formulation}[chapter]
\newtheorem{definition}{Definition}[chapter]
\newtheorem{lemma}[definition]{Lemma}
\newtheorem{theorem}[definition]{Theorem}
\newtheorem{corollary}[definition]{Corollary}
\theoremstyle{remark}
\newtheorem*{remark}{Remark}

%% Key points, contributions (C1, C2, ...), wide figures and tables, QR links: see environments.tex
%% =====================================================================
%% Environments for the text: key points, contributions, QR links, grids of panels, figure and table pages.
%% (Loaded by packages.tex.)
%% =====================================================================

%% ---------------------------------------------------------------------
%% Key points: a short list set apart between two thin rules, with a small heading, e.g. at the start
%% or end of a chapter. Tagged as an Aside, so a screen reader announces it as a side note.
%%   \begin{keypoints}            \begin{keypoints}[In short]
%%       \item ...                    \item ...
%%   \end{keypoints}              \end{keypoints}
%% ---------------------------------------------------------------------
\newcommand{\keypointsname}{Key points}
\ExplSyntaxOn
\NewDocumentEnvironment{keypoints}{ O{\keypointsname} }
  {
    \par \addvspace { \bigskipamount }
    \tag_struct_begin:n { tag = Aside }
    \hrule height 0.4pt \relax
    \nointerlineskip \vspace { 0.6\baselineskip }
    \noindent { \headingfont \bfseries \small \MakeUppercase {#1} } \par
    \begin{itemize}
  }
  {
    \end{itemize}
    \par \vspace { 0.4\baselineskip }
    \hrule height 0.4pt \relax
    \tag_struct_end:
    \par \addvspace { \bigskipamount }
  }
\ExplSyntaxOff

%% ---------------------------------------------------------------------
%% Contributions: a numbered list C1, C2, ... whose items can be referred to:
%%   \begin{contributions}
%%       \item An adjoint method for ... \label{c:adjoint}
%%   \end{contributions}
%%   ... as shown in \cref{c:adjoint}   gives "contribution C1"
%% ---------------------------------------------------------------------
\newlist{contributions}{enumerate}{1}
\setlist[contributions]{label=C\arabic*, font=\bfseries, leftmargin=2.4em} % (with tagging, LaTeX's own lists take enumitem's place and know fewer keys: no ref, leftmargin=* or widest)
\crefname{contributionsi}{contribution}{contributions}
\Crefname{contributionsi}{Contribution}{Contributions}

%% ---------------------------------------------------------------------
%% QR links: a QR code of a web address, with the address under it, in both editions, so the pages look
%% the same on paper and on screen; in the digital edition the code and the address are a link. For the
%% DOIs of data and code. Tagged as a figure with alternative text.
%%   \qrlink{https://doi.org/10.4121/00000000}     \qrlink[15mm]{...}   size (default 12 mm)
%%   \qrlink{...}[Data]                            a short label under it instead of the address
%% ---------------------------------------------------------------------
\usepackage[nolinks]{qrcode} % https://ctan.org/pkg/qrcode %% (its own links are switched off: ours below)
\ExplSyntaxOn
\NewDocumentCommand{\qrlink}{ O{12mm} m o }
  {
    \leavevmode
    \begin{tabular}[b]{@{}c@{}}
      \href{#2}
        {
          \tag_struct_begin:n { tag = Figure, alt = { QR~code~linking~to~#2 } }
          \tag_mc_begin:n { }
          \tag_suspend:n { qrcode }
          \color{black} \qrcode[height=#1]{#2} % black, not the link colour: the best contrast for scanning
          \tag_resume:n { qrcode }
          \tag_mc_end:
          \tag_struct_end:
        } \\
      \footnotesize \IfNoValueTF {#3} { \url{#2} } {#3}
    \end{tabular}
  }
\ExplSyntaxOff

%% ---------------------------------------------------------------------
%% Panels: a grid of plots of the same size, lined up in rows and columns, e.g. snapshots of a flow, each
%% with a short label under it and lettered (a), (b), ... automatically, and one shared colour bar under the
%% grid. In a figure, or in a figure page (below), where the grid is made to fit the page.
%%   \begin{figure}[tbp]
%%       \begin{panels}[columns=3]
%%           \panel{figures/snap-01.pdf}{$t = 0$}
%%           \panel[alt={...}, label={fig:snap-b}]{figures/snap-02.pdf}{$t = 0.35$}
%%           ...
%%           \colorbar{figures/colourbar.pdf}
%%       \end{panels}
%%       \caption{...}\label{fig:snapshots}
%%   \end{figure}
%% \cref{fig:snap-b} gives "Figure 2.9b", \subref{fig:snap-b} "(b)".
%% Options of panels:
%%   columns=3                  panels per row
%%   labels=below|inside|inside-right|none   the label under each panel (default), in its top left
%%                              (or right) corner on a white ground, or none
%%   letters=true|false         (a), (b), ... before the labels (default true)
%%   gap=2mm                    between the panels, across and down (natural: between the plot areas, 2.5mm)
%%   aspect=4/3                 width / height of the panels; default: that of the first image
%%   natural                    plots with axes: every file at its own size, not scaled, lined up on their
%%                              plot areas (files made by save_panels in figures/plots/example_figures.py: the same
%%                              plot area in all, the tick labels only in the left column and bottom row)
%%   xlabel={$x$}, ylabel={$u$} one axis title under and one beside the grid, centred on the plot areas;
%%   xlabel-sep=1pt, ylabel-sep=1.5mm   their distance from the tick labels
%%   pad=2mm                    natural: the margin round every plot area in the files (as in save_panels)
%% Options of \panel: alt={...} the alternative text for screen readers (default: "Panel b: " and the label),
%% label={...} for cross-references. \colorbar[alt={...}]{file}: an image at its own size (at most the width
%% of the grid), centred under it. Make the panels and the colour bar with the same colour scale.
%% (For a few sub-figures of different sizes, each with a caption of its own, use subfigure instead.)
%% ---------------------------------------------------------------------
\ExplSyntaxOn
\seq_new:N \g__thesis_panels_seq
\tl_new:N \g__thesis_panels_cbar_tl
\tl_new:N \g__thesis_panels_cbaralt_tl
\int_new:N \l__thesis_panels_columns_int
\int_new:N \l__thesis_panels_rows_int
\int_new:N \l__thesis_panels_count_int
\int_new:N \g__thesis_panels_last_int
\tl_new:N \l__thesis_panels_labels_tl
\bool_new:N \l__thesis_panels_letters_bool
\dim_new:N \l__thesis_panels_gap_dim
\tl_new:N \l__thesis_panels_aspect_tl
\dim_new:N \l__thesis_panels_wd_dim
\dim_new:N \l__thesis_panels_ht_dim
\dim_new:N \l__thesis_panels_labelht_dim
\dim_new:N \l__thesis_panels_grid_dim
\dim_new:N \l__thesis_panels_avail_dim
\dim_new:N \l__thesis_panels_maxht_dim      % the height the grid may take (set by figurepage; in a
                                            % figure: the text height less room for the caption)
\bool_new:N \l__thesis_panels_onpage_bool   % in a figure page: the colour bar goes to the foot
\box_new:N \l__thesis_panels_cbar_box
\box_new:N \l__thesis_panels_label_box
\fp_new:N \l__thesis_panels_aspect_fp
\box_new:N \l__thesis_panels_grid_box
\int_new:N \l__thesis_panels_col_int
\bool_new:N \l__thesis_panels_natural_bool
\tl_new:N \l__thesis_panels_xlabel_tl
\tl_new:N \l__thesis_panels_ylabel_tl
\dim_new:N \l__thesis_panels_left_dim       % natural: the extra width of the left column (y tick labels)
\dim_new:N \l__thesis_panels_bottom_dim     % natural: the extra height of the bottom row (x tick labels)
\dim_new:N \l__thesis_panels_pad_dim        % natural: the margin round every plot area (save_panels)
\dim_new:N \l__thesis_panels_inset_x_dim
\bool_new:N \l__thesis_panels_gapset_bool
\dim_new:N \l__thesis_panels_xsep_dim \dim_new:N \l__thesis_panels_ysep_dim
\dim_new:N \l__thesis_panels_padkey_dim \dim_set:Nn \l__thesis_panels_padkey_dim { 2mm }
\keys_define:nn { thesis / panels }
  {
    columns .int_set:N = \l__thesis_panels_columns_int ,
    columns .initial:n = 2 ,
    labels .choices:nn = { below , inside , inside-right , none } { \tl_set_eq:NN \l__thesis_panels_labels_tl \l_keys_choice_tl } ,
    labels .initial:n = below ,
    letters .bool_set:N = \l__thesis_panels_letters_bool ,
    letters .initial:n = true ,
    letters .default:n = true ,
    gap .code:n = \dim_set:Nn \l__thesis_panels_gap_dim {#1} \bool_set_true:N \l__thesis_panels_gapset_bool ,
    gap .initial:n = 2mm ,
    aspect .tl_set:N = \l__thesis_panels_aspect_tl ,
    natural .bool_set:N = \l__thesis_panels_natural_bool ,
    natural .default:n = true ,
    xlabel .tl_set:N = \l__thesis_panels_xlabel_tl ,
    ylabel .tl_set:N = \l__thesis_panels_ylabel_tl ,
    xlabel-sep .dim_set:N = \l__thesis_panels_xsep_dim ,
    xlabel-sep .initial:n = 1pt ,
    ylabel-sep .dim_set:N = \l__thesis_panels_ysep_dim ,
    ylabel-sep .initial:n = 1.5mm ,
    pad .code:n = \dim_set:Nn \l__thesis_panels_padkey_dim {#1} ,
  }
\keys_define:nn { thesis / colorbar } { alt .code:n = \tl_gset:Nn \g__thesis_panels_cbaralt_tl {#1} }
\keys_define:nn { thesis / panel }
  {
    alt .tl_set:N = \l__thesis_panel_alt_tl ,
    label .tl_set:N = \l__thesis_panel_label_tl ,
    data .tl_set:N = \l__thesis_panel_data_tl ,
  }

\NewDocumentEnvironment { panels } { O{} +b }
  {
    \par
    \group_begin:
    \bool_set_false:N \l__thesis_panels_gapset_bool
    \keys_set:nn { thesis / panels } {#1}
    % 1. collect the panels and the colour bar
    \seq_gclear:N \g__thesis_panels_seq
    \tl_gclear:N \g__thesis_panels_cbar_tl
    \RenewDocumentCommand \panel { O{} m m } { \seq_gput_right:Nn \g__thesis_panels_seq { {##1} {##2} {##3} } }
    \RenewDocumentCommand \colorbar { O{} m }
      {
        \tl_gclear:N \g__thesis_panels_cbaralt_tl
        \keys_set:nn { thesis / colorbar } {##1}
        \tl_gset:Nn \g__thesis_panels_cbar_tl {##2}
      }
    \hbox_set:Nn \l_tmpa_box { \tag_suspend:n { panels } #2 \tag_resume:n { panels } }
    \int_set:Nn \l__thesis_panels_count_int { \seq_count:N \g__thesis_panels_seq }
    \int_compare:nNnT { \l__thesis_panels_count_int } > { 0 } { \__thesis_panels_set: }
    \group_end:
    % the next figure page with continued goes on with the letters
    \int_gset:Nn \g__thesis_panels_last_int { \value { subfigure } }
  }
  { }
% \panel and \colorbar outside panels
\NewDocumentCommand \panel { O{} m m } { \msg_error:nn { thesis } { panel-outside } }
\NewDocumentCommand \colorbar { O{} m } { \msg_error:nn { thesis } { panel-outside } }
\msg_new:nnn { thesis } { panel-outside } { \token_to_str:N \panel\ and~ \token_to_str:N \colorbar\ belong~ inside~ a~ panels~ environment. }

% the image of a panel, measured without tagging
\cs_new_protected:Npn \__thesis_panels_measure:Nn #1#2
  { \hbox_set:Nn #1 { \tag_suspend:n { panels } \includegraphics {#2} \tag_resume:n { panels } } }

\cs_new_protected:Npn \__thesis_panels_set:
  {
    \int_set:Nn \l__thesis_panels_rows_int
      { \int_div_truncate:nn { \l__thesis_panels_count_int + \l__thesis_panels_columns_int - 1 } { \l__thesis_panels_columns_int } }
    \bool_set_eq:NN \l__thesis_subcaption_letter_bool \l__thesis_panels_letters_bool
    \skip_set:Nn \abovecaptionskip { 3pt } % the labels close to their panels
    % the colour bar, measured
    \tl_if_empty:NF \g__thesis_panels_cbar_tl
      { \__thesis_panels_measure:Nn \l__thesis_panels_cbar_box { \g__thesis_panels_cbar_tl } }
    \bool_if:NTF \l__thesis_panels_natural_bool
      { \__thesis_panels_size_natural: }
      { \__thesis_panels_size_scaled: }
    % the grid, row by row, in a box as wide as the grid
    \vbox_set:Nn \l__thesis_panels_grid_box
      {
        \dim_set:Nn \hsize { \l__thesis_panels_grid_dim } \dim_set_eq:NN \linewidth \hsize
        \skip_zero:N \parskip \skip_set:Nn \parfillskip { 0pt plus 1fil }
        \int_zero:N \l_tmpa_int
        \seq_map_inline:Nn \g__thesis_panels_seq
          {
            \int_set:Nn \l__thesis_panels_col_int { \int_mod:nn { \l_tmpa_int } { \l__thesis_panels_columns_int } }
            \int_compare:nNnTF { \l__thesis_panels_col_int } = { 0 }
              {
                \int_compare:nNnF { \l_tmpa_int } = { 0 } { \par \nointerlineskip \vspace { \l__thesis_panels_gap_dim } }
                \noindent
              }
              { \skip_horizontal:N \l__thesis_panels_gap_dim }
            \__thesis_panels_one:nnn ##1
            \int_incr:N \l_tmpa_int
          }
        \par
      }
    % the grid with the shared axis titles: y beside it, x under it, centred on the plot areas
    \tl_if_empty:NF \l__thesis_panels_ylabel_tl
      {
        \hbox_set:Nn \l_tmpa_box
          {
            \rotatebox { 90 }
              {
                \hbox_to_wd:nn
                  {
                    \box_ht:N \l__thesis_panels_grid_box
                    - \l__thesis_panels_bottom_dim - 2\l__thesis_panels_pad_dim
                  }
                  { \hss \small \l__thesis_panels_ylabel_tl \hss }
              }
          }
      }
    \noindent
    \hbox_to_wd:nn { \linewidth }
      {
        \hss
        \tl_if_empty:NF \l__thesis_panels_ylabel_tl
          {
            \box_move_up:nn { \l__thesis_panels_bottom_dim + \l__thesis_panels_pad_dim } { \box_use:N \l_tmpa_box }
            \skip_horizontal:N \l__thesis_panels_ysep_dim
          }
        \box_use_drop:N \l__thesis_panels_grid_box
        \hss
      }
    \par
    \tl_if_empty:NF \l__thesis_panels_xlabel_tl
      {
        \nointerlineskip \vspace { \l__thesis_panels_xsep_dim }
        \noindent
        \hbox_to_wd:nn { \linewidth }
          {
            \hss
            \tl_if_empty:NF \l__thesis_panels_ylabel_tl { \skip_horizontal:n { \box_wd:N \l_tmpa_box + \l__thesis_panels_ysep_dim } }
            \skip_horizontal:n { \l__thesis_panels_left_dim + \l__thesis_panels_pad_dim }
            \hbox_to_wd:nn
              { \l__thesis_panels_grid_dim - \l__thesis_panels_left_dim - 2\l__thesis_panels_pad_dim }
              { \hss \small \l__thesis_panels_xlabel_tl \hss }
            \skip_horizontal:n { \l__thesis_panels_pad_dim }
            \hss
          }
        \par
      }
    % the colour bar: under the grid, or at the foot of a figure page
    \tl_if_empty:NF \g__thesis_panels_cbar_tl
      {
        \bool_if:NTF \l__thesis_panels_onpage_bool { \vfill } { \vspace { 2\l__thesis_panels_gap_dim } }
        \noindent \hbox_to_wd:nn { \linewidth }
          {
            \hss
            \includegraphics
              [ alt = { \tl_if_empty:NTF \g__thesis_panels_cbaralt_tl { Colour~bar } { \g__thesis_panels_cbaralt_tl } } ,
                width = { \dim_min:nn { \box_wd:N \l__thesis_panels_cbar_box } { \linewidth } } ]
              { \g__thesis_panels_cbar_tl }
            \hss
          }
        \par
        \bool_if:NT \l__thesis_panels_onpage_bool { \bool_gset_true:N \g__thesis_page_filled_bool }
      }
  }
% scaled: every panel the same size, from the width, made smaller if the grid would be too high
\cs_new_protected:Npn \__thesis_panels_size_scaled:
  {
    \dim_zero:N \l__thesis_panels_left_dim \dim_zero:N \l__thesis_panels_bottom_dim
    \dim_zero:N \l__thesis_panels_pad_dim
    \tl_if_empty:NTF \l__thesis_panels_aspect_tl
      {
        \seq_get_left:NN \g__thesis_panels_seq \l_tmpa_tl
        \__thesis_panels_measure:Nn \l_tmpb_box { \exp_last_unbraced:No \use_ii:nnn \l_tmpa_tl }
        \fp_set:Nn \l__thesis_panels_aspect_fp { \box_wd:N \l_tmpb_box / \box_ht:N \l_tmpb_box }
      }
      { \fp_set:Nn \l__thesis_panels_aspect_fp { \l__thesis_panels_aspect_tl } }
    \hbox_set:Nn \l__thesis_panels_label_box { \small \strut }
    \dim_set:Nn \l__thesis_panels_labelht_dim
      {
        \str_if_eq:VnTF \l__thesis_panels_labels_tl { below }
          { \box_ht:N \l__thesis_panels_label_box + \box_dp:N \l__thesis_panels_label_box + \abovecaptionskip + \belowcaptionskip + 2pt } { 0pt }
      }
    \bool_if:NF \l__thesis_panels_onpage_bool
      { \dim_set:Nn \l__thesis_panels_maxht_dim { \textheight - 4\baselineskip } }
    \dim_set:Nn \l__thesis_panels_avail_dim
      {
        ( \linewidth
          - \tl_if_empty:NTF \l__thesis_panels_ylabel_tl { 0pt } { \baselineskip + \l__thesis_panels_ysep_dim }
          - \l__thesis_panels_gap_dim * ( \l__thesis_panels_columns_int - 1 ) ) / \l__thesis_panels_columns_int
      }
    \dim_set_eq:NN \l__thesis_panels_wd_dim \l__thesis_panels_avail_dim
    \dim_set:Nn \l__thesis_panels_ht_dim { \fp_to_dim:n { \l__thesis_panels_wd_dim / \l__thesis_panels_aspect_fp } }
    \dim_set:Nn \l__thesis_panels_avail_dim
      {
        \l__thesis_panels_maxht_dim
        - \l__thesis_panels_gap_dim * ( \l__thesis_panels_rows_int - 1 )
        - \l__thesis_panels_labelht_dim * \l__thesis_panels_rows_int
        - \tl_if_empty:NTF \l__thesis_panels_xlabel_tl { 0pt } { \baselineskip + \l__thesis_panels_xsep_dim }
        - \tl_if_empty:NTF \g__thesis_panels_cbar_tl { 0pt }
            { \box_ht:N \l__thesis_panels_cbar_box + \box_dp:N \l__thesis_panels_cbar_box + 2\l__thesis_panels_gap_dim }
      }
    \dim_compare:nNnT { \l__thesis_panels_ht_dim * \l__thesis_panels_rows_int } > { \l__thesis_panels_avail_dim }
      {
        \dim_set:Nn \l__thesis_panels_ht_dim { \l__thesis_panels_avail_dim / \l__thesis_panels_rows_int }
        \dim_set:Nn \l__thesis_panels_wd_dim { \fp_to_dim:n { \l__thesis_panels_ht_dim * \l__thesis_panels_aspect_fp } }
      }
    \dim_set:Nn \l__thesis_panels_grid_dim
      { \l__thesis_panels_wd_dim * \l__thesis_panels_columns_int + \l__thesis_panels_gap_dim * ( \l__thesis_panels_columns_int - 1 ) }
  }
% natural: every file at its own size (made by save_panels in Python: the same plot area in all, the tick
% labels in the files of the left column and the bottom row); their extra width and height, from the sizes
\cs_new_protected:Npn \__thesis_panels_size_natural:
  {
    \dim_set_eq:NN \l__thesis_panels_pad_dim \l__thesis_panels_padkey_dim
    % gap: between the plot areas (2.5 mm unless given); the files, with their margins, overlap by the rest
    \bool_if:NF \l__thesis_panels_gapset_bool { \dim_set:Nn \l__thesis_panels_gap_dim { 2.5mm } }
    \dim_set:Nn \l__thesis_panels_gap_dim { \l__thesis_panels_gap_dim - 2\l__thesis_panels_pad_dim }
    \seq_get_left:NN \g__thesis_panels_seq \l_tmpa_tl
    \__thesis_panels_measure:Nn \l_tmpa_box { \exp_last_unbraced:No \use_ii:nnn \l_tmpa_tl }
    \dim_zero:N \l__thesis_panels_left_dim \dim_zero:N \l__thesis_panels_bottom_dim
    \dim_set:Nn \l__thesis_panels_grid_dim { \box_wd:N \l_tmpa_box }
    \int_compare:nNnT { \l__thesis_panels_columns_int } > { 1 }
      {
        \tl_set:Ne \l_tmpb_tl { \seq_item:Nn \g__thesis_panels_seq { 2 } }
        \__thesis_panels_measure:Nn \l_tmpb_box { \exp_last_unbraced:No \use_ii:nnn \l_tmpb_tl }
        \dim_set:Nn \l__thesis_panels_left_dim { \box_wd:N \l_tmpa_box - \box_wd:N \l_tmpb_box }
        \dim_set:Nn \l__thesis_panels_grid_dim
          {
            \box_wd:N \l_tmpa_box
            + ( \box_wd:N \l_tmpb_box + \l__thesis_panels_gap_dim ) * ( \l__thesis_panels_columns_int - 1 )
          }
      }
    \int_compare:nNnT { \l__thesis_panels_rows_int } > { 1 }
      {
        \tl_set:Ne \l_tmpb_tl
          { \seq_item:Nn \g__thesis_panels_seq { ( \l__thesis_panels_rows_int - 1 ) * \l__thesis_panels_columns_int + 1 } }
        \__thesis_panels_measure:Nn \l_tmpb_box { \exp_last_unbraced:No \use_ii:nnn \l_tmpb_tl }
        \dim_set:Nn \l__thesis_panels_bottom_dim { \box_ht:N \l_tmpb_box - \box_ht:N \l_tmpa_box }
      }
  }
% a label in the top left (or right) corner of the plot area, on a white ground: label, left or right
\cs_new_protected:Npn \__thesis_panels_inside:nn #1#2
  {
    \hbox_set:Nn \l_tmpa_box
      {
        \small \bool_if:NT \l__thesis_panels_letters_bool
          { ( \int_eval:n { \value { subfigure } + 1 } ) \nobreakspace } #1
      }
    \vbox_set:Nn \l_tmpa_box
      {
        % as wide as the label (at most the panel), set from the left
        \dim_set:Nn \hsize { \dim_min:nn { \box_wd:N \l_tmpa_box + 2pt } { \l__thesis_panels_wd_dim } }
        \dim_set:Nn \linewidth { \hsize }
        \skip_set:Nn \abovecaptionskip { 1pt } \skip_set:Nn \belowcaptionskip { 0pt }
        \legacy_if_set_false:n { thesis@capcentre }
        \caption {#1}
        \__thesis_panels_label:
      }
    \hbox_set:Nn \l_tmpb_box
      {
        \color { white } \vrule width \box_wd:N \l_tmpa_box height \box_ht:N \l_tmpa_box depth \box_dp:N \l_tmpa_box
        \skip_horizontal:n { - \box_wd:N \l_tmpa_box } \color { black } \box_use:N \l_tmpa_box
      }
    \str_if_eq:nnT {#2} { right }
      { \dim_set:Nn \l__thesis_panels_inset_x_dim { \l__thesis_panels_wd_dim - \l__thesis_panels_pad_dim - \box_wd:N \l_tmpb_box } }
    \par \nointerlineskip
    \vspace* { \l__thesis_panels_pad_dim - \l__thesis_panels_ht_dim }
    \hbox:n { \skip_horizontal:N \l__thesis_panels_inset_x_dim \box_use:N \l_tmpb_box }
    \nointerlineskip
    \vspace* { \dim_eval:n { \l__thesis_panels_ht_dim - \l__thesis_panels_pad_dim - \box_ht:N \l_tmpb_box - \box_dp:N \l_tmpb_box } }
  }
% one panel: options, file, label
\cs_new_protected:Npn \__thesis_panels_label:
  { \tl_if_empty:NF \l__thesis_panel_label_tl { \exp_args:NV \label \l__thesis_panel_label_tl } }
\cs_new_protected:Npn \__thesis_panels_one:nnn #1#2#3
  {
    \group_begin:
    \tl_clear:N \l__thesis_panel_alt_tl \tl_clear:N \l__thesis_panel_label_tl \tl_clear:N \l__thesis_panel_data_tl
    \keys_set:nn { thesis / panel } {#1}
    \tl_if_empty:NT \l__thesis_panel_alt_tl % (default: "Panel b: t = 0.35", from the label without its $)
      {
        \tl_set:Ne \l__thesis_panel_alt_tl { \text_purify:n {#3} }
        \tl_remove_all:Nn \l__thesis_panel_alt_tl { $ }
        \tl_set:Ne \l__thesis_panel_alt_tl
          {
            Panel \bool_if:NT \l__thesis_panels_letters_bool { ~ \int_to_alph:n { \value { subfigure } + 1 } }
            \tl_if_blank:VF \l__thesis_panel_alt_tl { : ~ \exp_not:V \l__thesis_panel_alt_tl }
          }
      }
    % the size of this panel; the corner of its plot area (natural: inside the tick labels and the pad)
    \dim_zero:N \l__thesis_panels_inset_x_dim
    \bool_if:NT \l__thesis_panels_natural_bool
      {
        \__thesis_panels_measure:Nn \l_tmpb_box {#2}
        \dim_set:Nn \l__thesis_panels_wd_dim { \box_wd:N \l_tmpb_box }
        \dim_set:Nn \l__thesis_panels_ht_dim { \box_ht:N \l_tmpb_box }
        \dim_set:Nn \l__thesis_panels_inset_x_dim
          { \l__thesis_panels_pad_dim + \int_compare:nNnTF { \l__thesis_panels_col_int } = { 0 } { \l__thesis_panels_left_dim } { 0pt } }
      }
    \begin{subfigure}[t]{\l__thesis_panels_wd_dim}
      \includegraphics [ alt = { \l__thesis_panel_alt_tl } , width = \l__thesis_panels_wd_dim , height = \l__thesis_panels_ht_dim ] {#2}
      \tl_if_empty:NF \l__thesis_panel_data_tl
        { \use:e { \exp_not:N \figuredata [ Data~of~\exp_not:V \l__thesis_panel_alt_tl ] { \exp_not:V \l__thesis_panel_data_tl } } }
      \str_case:Vn \l__thesis_panels_labels_tl
        {
          { below } { \caption {#3} }
          { inside } { \__thesis_panels_inside:nn {#3} { left } }
          { inside-right } { \__thesis_panels_inside:nn {#3} { right } }
          { none } { \refstepcounter { subfigure } }
        }
      \str_if_in:NnF \l__thesis_panels_labels_tl { inside } { \__thesis_panels_label: }
    \end{subfigure}
    \group_end:
  }
\ExplSyntaxOff

%% ---------------------------------------------------------------------
%% Figure and table pages: a figure or a table on a page of its own, as a float (the text goes on around
%% it), e.g. a grid of panels that needs the whole page.
%%   \begin{figurepage}                        \begin{tablepage}[landscape]
%%       \begin{panels}[columns=2] ... \end{panels}     \caption{...}\label{tab:...}
%%       \caption{...}\label{fig:...}                   \begin{tabular}...\end{tabular}
%%   \end{figurepage}                          \end{tablepage}
%% or a single large figure: \includegraphics[width=\linewidth, height=\areaheight, keepaspectratio]{...}
%% Options:
%%   landscape     turned sideways, like sidewaysfigure (shown upright on screen in the digital edition)
%%   wide          also the outer margin, up to 10 mm from the edge of the page (the caption keeps the
%%                 width of the text)
%%   continued     (figurepage) goes on from the previous figure page: the same number, "(continued)" in
%%                 the caption, no second entry in the list of figures, panel letters carried on
%%   head=false    without the running head and page number (default: with them, as on every page)
%%   align=centre  the content in the middle of the page (default: at the top); the caption, and the colour
%%                 bar of a panel grid, stay at the foot (a table's caption at the top)
%% \linewidth is the width of the area, \areaheight its height less room for a three-line caption. A long
%% table that runs over several pages is an xltabular, not a table page.
%% ---------------------------------------------------------------------
\newlength{\areaheight}
\ExplSyntaxOn
\bool_new:N \l__thesis_page_landscape_bool
\bool_new:N \l__thesis_page_wide_bool
\bool_new:N \l__thesis_page_continued_bool
\bool_new:N \l__thesis_page_head_bool
\bool_new:N \l__thesis_page_centre_bool
\bool_new:N \g__thesis_page_filled_bool
\dim_new:N \l__thesis_page_extra_dim
\dim_new:N \l__thesis_page_wd_dim
\dim_new:N \l__thesis_page_ht_dim
\keys_define:nn { thesis / page }
  {
    landscape .bool_set:N = \l__thesis_page_landscape_bool ,
    landscape .default:n = true ,
    wide .bool_set:N = \l__thesis_page_wide_bool ,
    wide .default:n = true ,
    continued .bool_set:N = \l__thesis_page_continued_bool ,
    continued .default:n = true ,
    head .bool_set:N = \l__thesis_page_head_bool ,
    head .initial:n = true ,
    head .default:n = true ,
    align .choices:nn = { top , centre , center }
      { \bool_set:Nn \l__thesis_page_centre_bool { ! \str_if_eq_p:Vn \l_keys_choice_tl { top } } } ,
  }

%% Each is a float of its own (placement p) that fills the page, so LaTeX keeps it in order with the other
%% floats. Whether its page is a left or a right one, and which page goes without running head, is read
%% from the previous run (it decides on which side wide goes).
\int_new:N \g__thesis_page_int
\cs_new_protected:Npn \__thesis_page_mark:
  {
    \iow_shipout_x:cx { @auxout }
      { \exp_not:c { thesis@pagefloat } { \int_use:N \g__thesis_page_int } { \exp_not:N \the \exp_not:N \value { page } } }
    \bool_if:NF \l__thesis_page_head_bool
      { \iow_shipout_x:cx { @auxout } { \exp_not:c { thesis@nohead } { \exp_not:N \the \exp_not:N \value { page } } } }
  }
\cs_new_protected:cpn { thesis@pagefloat } #1#2 { \cs_gset:cpn { thesis@pagefloat@#1 } {#2} }
\cs_new_protected:cpn { thesis@nohead } #1 { \cs_gset:cpn { thesis@nohead@#1 } { } }
\cs_new:Npn \__thesis_page_page: { \cs_if_exist_use:cF { thesis@pagefloat@ \int_use:N \g__thesis_page_int } { 1 } }
% a page marked without head: no running head and foot
\AddToHook { cmd/@outputpage/before }
  {
    \cs_if_exist:cT { thesis@nohead@ \int_use:N \value { page } }
      { \clist_map_inline:nn { @oddhead , @evenhead , @oddfoot , @evenfoot } { \cs_set_eq:cN {##1} \prg_do_nothing: } }
  }

\NewDocumentEnvironment { figurepage } { O{} +b } { \__thesis_page:nnn { figure } {#1} {#2} } { }
\NewDocumentEnvironment { tablepage } { O{} +b } { \__thesis_page:nnn { table } {#1} {#2} } { }

\cs_new_protected:Npn \__thesis_page:nnn #1#2#3 % figure or table, options, body
  {
    \group_begin:
    \int_gincr:N \g__thesis_page_int
    \keys_set:nn { thesis / page } {#2}
    \dim_set:Nn \l__thesis_page_extra_dim
      { \bool_if:NTF \l__thesis_page_wide_bool { \foremargin - 10mm } { 0pt } }
    % the area: width and height (turned: the text height wide)
    \bool_if:NTF \l__thesis_page_landscape_bool
      {
        \dim_set:Nn \l__thesis_page_wd_dim { \textheight }
        \dim_set:Nn \l__thesis_page_ht_dim { \textwidth + \l__thesis_page_extra_dim }
      }
      {
        \dim_set:Nn \l__thesis_page_wd_dim { \textwidth + \l__thesis_page_extra_dim }
        \dim_set:Nn \l__thesis_page_ht_dim { \textheight }
      }
    \dim_set:Nn \areaheight { \l__thesis_page_ht_dim - 4\baselineskip }
    \begin {#1} [p]
      \__thesis_page_mark:
      \bool_if:NTF \l__thesis_page_landscape_bool
        {
          \thesisrotatepage
          \hbox_to_wd:nn { \textwidth }
            {
              \int_if_even:nT { \__thesis_page_page: } { \skip_horizontal:n { - \l__thesis_page_extra_dim } }
              \rotatebox { 90 } { \__thesis_page_area:nn {#1} {#3} }
              \hss
            }
        }
        {
          \int_if_even:nT { \__thesis_page_page: } { \tex_moveleft:D \l__thesis_page_extra_dim }
          \__thesis_page_area:nn {#1} {#3}
        }
    \end {#1}
    \group_end:
  }
% the area of the page (a minipage would start a paragraph in the tagging, inside a box)
\cs_new_protected:Npn \__thesis_page_area:nn #1#2
  {
    \vbox_to_ht:nn { \l__thesis_page_ht_dim }
      {
        \dim_set:Nn \hsize { \l__thesis_page_wd_dim } \dim_set:Nn \linewidth { \hsize }
        \dim_set:Nn \textwidth { \hsize } \dim_set:Nn \columnwidth { \hsize }
        \dim_set_eq:NN \l__thesis_panels_maxht_dim \areaheight
        \bool_set_true:N \l__thesis_panels_onpage_bool
        \bool_gset_false:N \g__thesis_page_filled_bool
        \bool_if:NT \l__thesis_page_centre_bool { \vfil }
        \str_if_eq:nnTF {#1} { figure }
          {
            \bool_if:NT \l__thesis_page_continued_bool { \__thesis_page_continued: }
            % the caption at the foot of the page, with the text width
            \cs_set_eq:NN \__thesis_page_caption:w \caption
            \RenewDocumentCommand \caption { o m }
              {
                \bool_if:NF \g__thesis_page_filled_bool { \vfill }
                \__thesis_page_narrow:n
                  {
                    \bool_if:NTF \l__thesis_page_continued_bool
                      { \__thesis_page_caption:w [ ] { (continued) \tl_if_blank:nF {##2} { ~ ##2 } } }
                      { \IfValueTF {##1} { \__thesis_page_caption:w [##1] {##2} } { \__thesis_page_caption:w {##2} } }
                  }
                \bool_gset_true:N \g__thesis_page_captioned_bool
              }
            \bool_gset_false:N \g__thesis_page_captioned_bool
            #2
            \bool_if:NT \l__thesis_page_continued_bool
              { \bool_if:NF \g__thesis_page_captioned_bool { \caption { } } }
            \bool_if:NF \g__thesis_page_filled_bool { \vfill }
          }
          { #2 \vfill }
        \bool_if:NT \l__thesis_page_centre_bool { \vfil }
      }
  }
\bool_new:N \g__thesis_page_captioned_bool
% the caption of a wide figure keeps the width of the text, on the text's side of the area
\cs_new_protected:Npn \__thesis_page_narrow:n #1
  {
    \dim_compare:nNnTF { \l__thesis_page_extra_dim } = { 0pt }
      {#1}
      {
        \group_begin:
          \dim_set:Nn \hsize { \l__thesis_page_wd_dim - \l__thesis_page_extra_dim }
          \dim_set_eq:NN \linewidth \hsize
          \int_if_even:nT { \__thesis_page_page: } { \skip_set:Nn \leftskip { \l__thesis_page_extra_dim } }
          #1
        \group_end:
      }
  }
% continued: the number of the previous figure, its panel letters carried on, not in the list of figures
\cs_new_protected:Npn \__thesis_page_continued:
  {
    \addtocounter { figure } { -1 }
    \setcounter { subfigure } { \g__thesis_panels_last_int }
    \cs_set_eq:NN \addcontentsline \use_none:nnn
  }
\ExplSyntaxOff

%% What screen readers get: words in maths, table headers, alternative text for formulas, acronyms, units,
%% languages (accessibility.tex)
%% =====================================================================
%% Accessibility: what screen readers get in the tagged PDF, beyond LaTeX's own tagging. Nothing here
%% changes the look of the pages. Loaded at the end of packages.tex.
%%   1. Words in maths: operator names (sin, tr, arg max) and \mathword{Re} are read as one word
%%   2. Tables: the first row of a table is its header (so each cell is read with its column heading)
%%   3. Formulas with an alternative text of your own: \formulaalt{...}
%%   4. Acronyms: the short form carries its expansion ("BEM" is read as "blade element momentum")
%%   5. Units read as words: \qty{10}{\metre\per\second} is read as "10 metres per second"
%%   6. Other languages: text in an otherlanguage environment is marked with its language (Dutch)
%% =====================================================================
\makeatletter
\ExplSyntaxOn

%% ---------------------------------------------------------------------
%% 1. Words in maths. In the MathML every letter of "sin" or "Re" is a variable of its own, which screen
%% readers spell out ("s i n"). Operator names (\sin, \max, \DeclareMathOperator, \operatorname) and
%% \mathword{...} get the whole word as one identifier. \mathword{Re} looks like \mathrm{Re}: use it for
%% names of more than one letter, such as dimensionless numbers (commands.tex).
%% ---------------------------------------------------------------------
\text_declare_purify_equivalent:Nn \vphantom { \use_none:n }
\text_declare_purify_equivalent:Nn \, { ~ }
\text_declare_purify_equivalent:Nn \; { ~ }
\cs_new:Npn \__thesis_mathword:n #1 { \lua_escape:e { \text_purify:n {#1} } }
\NewDocumentCommand \mathword { m } { \mathrm {#1} }
% inside: no structure elements of its own (a \vphantom in \max would leave one without a parent, as the
% word replaces what is inside)
\cs_new_eq:NN \__thesis_mathword_nostruct: \thesis_luamml_no_structure:
% (luamml, which makes the MathML, is loaded at the start of the document)
\AddToHook { begindocument }
  {
    \cs_if_exist:NT \luamml_annotate:nen
      {
        % amsmath's \qopname, which all operator names go through, with the name as one identifier
        \cs_set_protected:Npn \qopname #1#2#3
          {
            \luamml_annotate:nen { 1 }
              { nucleus = true, mathml = { [0] = 'mi', '\__thesis_mathword:n {#3}' } }
              { \mathop { \__thesis_mathword_nostruct: #1 \kern \c_zero_dim \use:c { operator@font } #3 } \use:c { n#2limits@ } }
          }
        \RenewDocumentCommand \mathword { m }
          {
            \luamml_annotate:nen { 1 }
              { nucleus = true, mathml = { [0] = 'mi', '\__thesis_mathword:n {#1}' } }
              { \mathord { \__thesis_mathword_nostruct: \mathrm {#1} } }
          }
      }
  }

%% \layouttable before a tabular that only places text (the committee, the colophon, the curriculum
%% vitae): its rows and cells are then tagged as text blocks, not as a table, so a screen reader reads
%% them as text and checkers do not ask for header cells. It holds until the end of the group.
\NewDocumentCommand { \layouttable } { }
  { \tag_if_active:T { \tagpdfsetup { table/tagging = div } } }

\tag_if_active:T
  {
%% ---------------------------------------------------------------------
%% 2. Tables. The first row of every table in a table float (and of xltabular and longtable) is marked as
%% its header row. A table with a header of two rows: \tagpdfsetup{table/header-rows={1,2}} in the
%% table, before the tabular; none: \tagpdfsetup{table/header-rows={}}. Tables used only to place text:
%% \layouttable (above).
%% ---------------------------------------------------------------------
    \AddToHook { env/table/begin } { \tagpdfsetup { table/header-rows = 1 } }
    \AddToHook { env/sidewaystable/begin } { \tagpdfsetup { table/header-rows = 1 } }
    \AddToHook { env/xltabular/begin } { \tagpdfsetup { table/header-rows = 1 } }
    \AddToHook { env/longtable/begin } { \tagpdfsetup { table/header-rows = 1 } }

%% A \makecell (or \thead) is a small tabular inside a cell: set as a layout table, so that a header that
%% breaks over two lines stays one header cell with text, not a table inside a header cell.
    \AddToHook { begindocument/before }
      {
        \cs_if_exist:NT \mcell@@@tabular
          {
            \cs_new_eq:NN \__thesis_mcell_tabular:w \mcell@@@tabular
            \cs_set:Npn \mcell@@@tabular #1#2 \@nil #3
              { \group_begin: \layouttable \__thesis_mcell_tabular:w {#1} #2 \@nil {#3} \group_end: }
          }
      }

%% Links: each link annotation also gets a short description (/Contents: "Go to destination ...", or the
%% web address). PDF/UA-2 does not ask for it, but checkers built on PDF/UA-1 (PAC, Acrobat) do, and
%% LaTeX adds it by itself only for PDF/UA-1.
    \AddToHook { begindocument }
      {
        \clist_map_inline:nn { GoTo , GoToR , URI }
          {
            \socket_if_exist:nT { hyp/link/#1/Contents }
              { \socket_assign_plug:nn { hyp/link/#1/Contents } { default } }
          }
      }

%% ---------------------------------------------------------------------
%% 3. Formulas: an alternative text of your own, for a formula whose MathML reads badly (a large matrix,
%% a diagram drawn in maths). Write it right before the formula:
%%   \formulaalt{The stiffness matrix: a block matrix of four blocks, ...}
%%   \begin{equation} ... \end{equation}
%% Screen readers then read that text instead of the formula.
%% ---------------------------------------------------------------------
    \tl_new:N \g__thesis_formula_alt_tl
    \NewDocumentCommand \formulaalt { m } { \tl_gset:Nn \g__thesis_formula_alt_tl {#1} }
    \NewTaggingSocketPlug { math/content } { thesis }
      {
        % as LaTeX's alt+source, which the PDF/UA-2 set-up uses
        \tl_set:Ne \l__math_content_alt_tl { \l__math_content_template_tl }
        \tl_set:Ne \l__math_content_AF_source_tl { \l__math_texsource_template_tl }
        \tl_set:Nn \l__math_content_AF_mathml_tl { }
        \tl_set:Nn \l__math_content_actual_tl { }
        \tl_if_empty:NF \g__thesis_formula_alt_tl
          {
            \tl_set:Ne \l__math_content_alt_tl { \text_purify:n { \g__thesis_formula_alt_tl } }
            \bool_set_true:N \l__tag_math_alt_bool
            \tl_gclear:N \g__thesis_formula_alt_tl
          }
      }
    \AssignTaggingSocketPlug { math/content } { thesis }

%% ---------------------------------------------------------------------
%% 4. Acronyms: \gls{bem} after its first use gives "BEM", marked with its expansion (/E), so a screen
%% reader can say "blade element momentum theory". (formatNomenclature.tex calls \thesisexpansion.)
%% ---------------------------------------------------------------------
    \cs_gset_protected:Npn \thesisexpansion #1#2 % expansion, text
      {
        \mode_if_horizontal:F { \leavevmode }
        \tag_mc_end_push:
        \use:e { \exp_not:N \tag_struct_begin:n { tag = Span , E = { \text_purify:n {#1} } } }
        \tag_mc_begin:n { }
        #2
        \tag_mc_end:
        \tag_struct_end:
        \tag_mc_begin_pop:n { }
      }

%% ---------------------------------------------------------------------
%% 5. Units read as words: \unit{\metre\per\second} carries the expansion "metres per second", and
%% \qty{10}{\metre\per\second} "10 metres per second" (singular after 1). Only units written with
%% siunitx's macros (\kilo\metre, \per, \squared, ...) are spelt out; a unit written as text, such as
%% \unit{m/s}, is left as it is. Units of your own: \thesisunitword{\furlong}{furlong}.
%% ---------------------------------------------------------------------
    \prop_new:N \g__thesis_unit_words_prop
    \prop_new:N \g__thesis_unit_prefix_prop
    \cs_gset_protected:Npn \thesisunitword #1#2
      { \prop_gput:Nen \g__thesis_unit_words_prop { \cs_to_str:N #1 } {#2} }
    \cs_generate_variant:Nn \prop_gput:Nnn { Nen }
    \cs_gset_protected:Npn \__thesis_unit_pair:w #1 / #2 \q_stop #3 { \prop_gput:Nnn #3 {#1} {#2} }
    \clist_map_inline:nn
      {
        metre/metre, meter/metre, second/second, gram/gram, kilogram/kilogram, kelvin/kelvin,
        ampere/ampere, mole/mole, candela/candela, newton/newton, pascal/pascal, joule/joule,
        watt/watt, hertz/hertz, volt/volt, ohm/ohm, coulomb/coulomb, farad/farad, henry/henry,
        tesla/tesla, weber/weber, siemens/siemens, radian/radian, steradian/steradian,
        degreeCelsius/degree~Celsius, celsius/degree~Celsius, degree/degree, arcminute/arcminute,
        arcsecond/arcsecond, litre/litre, liter/litre, hour/hour, minute/minute, day/day,
        tonne/tonne, bar/bar, percent/percent, electronvolt/electronvolt, decibel/decibel,
        neper/neper, bel/bel, hectare/hectare, astronomicalunit/astronomical~unit,
        dalton/dalton, knot/knot, nauticalmile/nautical~mile, angstrom/angstrom, lux/lux,
        lumen/lumen, becquerel/becquerel, gray/gray, sievert/sievert, katal/katal
      }
      { \__thesis_unit_pair:w #1 \q_stop \g__thesis_unit_words_prop }
    \clist_map_inline:nn
      {
        quecto/quecto, ronto/ronto, yocto/yocto, zepto/zepto, atto/atto, femto/femto, pico/pico,
        nano/nano, micro/micro, milli/milli, centi/centi, deci/deci, deca/deca, deka/deca,
        hecto/hecto, kilo/kilo, mega/mega, giga/giga, tera/tera, peta/peta, exa/exa,
        zetta/zetta, yotta/yotta, ronna/ronna, quetta/quetta
      }
      { \__thesis_unit_pair:w #1 \q_stop \g__thesis_unit_prefix_prop }
    \seq_new:N \l__thesis_unit_seq
    \tl_new:N \l__thesis_unit_tmp_tl
    \tl_new:N \l__thesis_unit_words_tl
    \tl_new:N \l__thesis_unit_number_tl
    \bool_new:N \l__thesis_unit_ok_bool
    \cs_generate_variant:Nn \prop_get:NnNTF { NV }
    \cs_generate_variant:Nn \regex_match:nnTF { nV }
    \cs_generate_variant:Nn \str_case:nnF { V }
    % each unit is kept as {square or cubic}{prefix and unit}{squared or cubed}; "per" as {}{}{}
    \cs_gset_protected:Npn \__thesis_unit_words:nn #1#2 % unit, plural
      {
        \seq_clear:N \l__thesis_unit_seq
        \bool_set_true:N \l__thesis_unit_ok_bool
        \tl_clear:N \l_tmpa_tl % a prefix waiting for its unit
        \tl_clear:N \l_tmpb_tl % square or cubic, waiting
        \tl_map_inline:nn {#1}
          {
            \token_if_cs:NTF ##1
              {
                \str_set:Ne \l_tmpa_str { \cs_to_str:N ##1 }
                \prop_get:NVNTF \g__thesis_unit_prefix_prop \l_tmpa_str \l__thesis_unit_tmp_tl
                  { \tl_set_eq:NN \l_tmpa_tl \l__thesis_unit_tmp_tl }
                  {
                    \prop_get:NVNTF \g__thesis_unit_words_prop \l_tmpa_str \l__thesis_unit_tmp_tl
                      {
                        \seq_put_right:Ne \l__thesis_unit_seq
                          { { \exp_not:V \l_tmpb_tl } { \exp_not:V \l_tmpa_tl \exp_not:V \l__thesis_unit_tmp_tl } { } }
                        \tl_clear:N \l_tmpa_tl \tl_clear:N \l_tmpb_tl
                      }
                      {
                        \str_case:VnF \l_tmpa_str
                          {
                            { per } { \seq_put_right:Nn \l__thesis_unit_seq { { } { } { } } }
                            { square } { \tl_set:Nn \l_tmpb_tl { square } }
                            { cubic } { \tl_set:Nn \l_tmpb_tl { cubic } }
                            { squared } { \__thesis_unit_suffix:n { squared } }
                            { cubed } { \__thesis_unit_suffix:n { cubed } }
                          }
                          { \bool_set_false:N \l__thesis_unit_ok_bool }
                      }
                  }
              }
              { \tl_if_blank:nF {##1} { \bool_set_false:N \l__thesis_unit_ok_bool } }
          }
        \tl_clear:N \l__thesis_unit_words_tl
        \bool_lazy_and:nnT { \l__thesis_unit_ok_bool } { ! \seq_if_empty_p:N \l__thesis_unit_seq }
          {
            % plural: the last unit before the first "per"
            \int_zero:N \l_tmpa_int \int_zero:N \l_tmpb_int
            \bool_if:nT {#2}
              {
                \seq_map_indexed_inline:Nn \l__thesis_unit_seq
                  {
                    \tl_if_blank:eTF { \use_ii:nnn ##2 }
                      { \int_compare:nNnT { \l_tmpb_int } = { 0 } { \int_set:Nn \l_tmpb_int {##1} } }
                      { \int_compare:nNnT { \l_tmpb_int } = { 0 } { \int_set:Nn \l_tmpa_int {##1} } }
                  }
              }
            \seq_map_indexed_inline:Nn \l__thesis_unit_seq
              { \__thesis_unit_join:nnnn {##1} ##2 }
          }
      }
    \cs_gset_protected:Npn \__thesis_unit_join:nnnn #1#2#3#4 % index, power, unit, suffix
      {
        \tl_put_right:Ne \l__thesis_unit_words_tl
          {
            \tl_if_empty:NF \l__thesis_unit_words_tl { ~ }
            \tl_if_blank:nTF {#3}
              { per }
              {
                \tl_if_blank:nF {#2} { #2 ~ }
                \int_compare:nNnTF {#1} = { \l_tmpa_int } { \__thesis_unit_s:n {#3} } {#3}
                \tl_if_blank:nF {#4} { ~ #4 }
              }
          }
      }
    \cs_gset_protected:Npn \__thesis_unit_suffix:n #1
      {
        \seq_pop_right:NNTF \l__thesis_unit_seq \l__thesis_unit_tmp_tl
          { \seq_put_right:Ne \l__thesis_unit_seq { \exp_last_unbraced:NV \__thesis_unit_setsuffix:nnnn \l__thesis_unit_tmp_tl {#1} } }
          { \bool_set_false:N \l__thesis_unit_ok_bool }
      }
    \cs_gset:Npn \__thesis_unit_setsuffix:nnnn #1#2#3#4 { { \exp_not:n {#1} } { \exp_not:n {#2} } { \exp_not:n {#4} } }
    % the plural of a unit word: an s, except for a few ("degree Celsius" -> "degrees Celsius")
    \cs_gset:Npn \__thesis_unit_s:n #1
      {
        \str_case:nnF {#1}
          {
            { hertz } {#1} { siemens } {#1} { lux } {#1} { percent } {#1}
            { degree~Celsius } { degrees~Celsius }
          }
          { #1 s }
      }
    % the wrappers: in text only (in maths the MathML has the units)
    \AddToHook { begindocument }
      {
        \NewCommandCopy \__thesis_qty:w \qty % (copies that keep working once \qty is redefined)
        \NewCommandCopy \__thesis_unit:w \unit
        \RenewDocumentCommand \qty { O{} m m }
          {
            \mode_if_math:TF { \__thesis_qty:w [#1] {#2} {#3} }
              {
                \tl_set:Ne \l__thesis_unit_number_tl { \tl_trim_spaces:n {#2} }
                \bool_set:Nn \l_tmpa_bool { ! \str_if_eq_p:Vn \l__thesis_unit_number_tl { 1 } }
                \__thesis_unit_words:nn {#3} { \l_tmpa_bool }
                % (only a plain number: one in another form is read as written)
                \regex_match:nVTF { \A[+\-]?[0-9]+([.,][0-9]+)?\Z } \l__thesis_unit_number_tl
                  {
                    \tl_if_empty:NTF \l__thesis_unit_words_tl
                      { \__thesis_qty:w [#1] {#2} {#3} }
                      { \thesisexpansion { #2 ~ \l__thesis_unit_words_tl } { \__thesis_qty:w [#1] {#2} {#3} } }
                  }
                  { \__thesis_qty:w [#1] {#2} {#3} }
              }
          }
        \RenewDocumentCommand \unit { O{} m }
          {
            \mode_if_math:TF { \__thesis_unit:w [#1] {#2} }
              {
                \__thesis_unit_words:nn {#2} { \c_true_bool }
                \tl_if_empty:NTF \l__thesis_unit_words_tl
                  { \__thesis_unit:w [#1] {#2} }
                  { \thesisexpansion { \l__thesis_unit_words_tl } { \__thesis_unit:w [#1] {#2} } }
              }
          }
      }

%% ---------------------------------------------------------------------
%% 6. Other languages: an otherlanguage environment (babel) is a division marked with its language, e.g.
%% nl for Dutch, and \foreignlanguage a span, so screen readers switch pronunciation. (A switch such as
%% \selectlanguage, or otherlanguage*, has no end to close a structure at: use otherlanguage instead.)
%% ---------------------------------------------------------------------
    \AddToHook { begindocument/before }
      {
        \cs_set_eq:NN \__thesis_otherlanguage:n \otherlanguage
        \cs_set_eq:NN \__thesis_endotherlanguage: \endotherlanguage
        \cs_set_protected:Npn \otherlanguage #1
          {
            \__thesis_otherlanguage:n {#1}
            \par
            \tag_struct_begin:n { tag = Div , lang = \localeinfo* { tag.bcp47 } }
          }
        \cs_set_protected:Npn \endotherlanguage
          { \par \tag_struct_end: \__thesis_endotherlanguage: }
        % a few words in another language, \foreignlanguage{dutch}{proefschrift}: a span with that language
        \NewCommandCopy \__thesis_foreignlanguage:w \foreignlanguage
        \RenewDocumentCommand \foreignlanguage { O{} m m }
          { \__thesis_foreignlanguage:w [#1] {#2} { \__thesis_langspan:n {#3} } }
      }
    \cs_new_protected:Npn \__thesis_langspan:n #1
      {
        \mode_if_horizontal:F { \leavevmode }
        \tag_mc_end_push:
        \use:e { \exp_not:N \tag_struct_begin:n { tag = Span , lang = \localeinfo* { tag.bcp47 } } }
        \tag_mc_begin:n { }
        #1
        \tag_mc_end:
        \tag_struct_end:
        \tag_mc_begin_pop:n { }
      }
  }
\cs_if_exist:NF \thesisexpansion { \cs_new_protected:Npn \thesisexpansion #1#2 {#2} }

\ExplSyntaxOff
\makeatother

%% =====================================================================
%% Files attached to the PDF (PDF/A-4f allows any file, as an "associated file" with a description).
%% PDF viewers list them in their attachments panel.
%%   1. A citation file: how to cite the thesis, as BibTeX (.bib) and RIS (.ris, for EndNote, Zotero,
%%      Mendeley), made from settings/thesisinfo.tex on every compile.
%%   2. \figuredata[description]{file}: the data behind a figure, after its \includegraphics (or the key
%%      data={file} of \panel). In the tagged PDF the file belongs to that image; screen readers and other
%%      tools can offer it.
%%   3. \supplementfile[description]{file}: any other file, attached to the document as a whole (input
%%      files, a spreadsheet, a short script).
%% Keep attachments small: they make the PDF larger, also in print (the printer ignores them, but they
%% are sent along). Large data sets belong in a repository such as 4TU.ResearchData, cited with their DOI.
%% =====================================================================
%% For the printer (PDF/X) and in the preview (a plain PDF): no attachments; the commands do nothing.
\newif\ifthesisnoattach
\ifdefined\thesispdfx \thesisnoattachtrue \fi
\ifdefined\thesisplainpdf \thesisnoattachtrue \fi
\ifthesisnoattach
\NewDocumentCommand \figuredata { O{} m } { }
\NewDocumentCommand \supplementfile { O{} m } { }
\expandafter\endinput
\fi
\ExplSyntaxOn

%% File types: PDF/A needs each attached file's type. LaTeX knows .csv, .txt, .pdf, .png, ...; a few more.
\prop_gput_from_keyval:Nn \g_pdffile_mimetypes_prop
  {
    .bib  = text/x-bibtex ,
    .ris  = application/x-research-info-systems ,
    .dat  = text/plain ,
    .tsv  = text/tab-separated-values ,
    .json = application/json ,
    .md   = text/markdown ,
    .py   = text/x-python ,
    .m    = text/plain ,
    .jl   = text/plain ,
    .zip  = application/zip ,
    .h5   = application/x-hdf5 ,
    .hdf5 = application/x-hdf5 ,
    .nc   = application/x-netcdf ,
    .npy  = application/octet-stream ,
    .npz  = application/zip ,
    .xlsx = application/vnd.openxmlformats-officedocument.spreadsheetml.sheet ,
    .ods  = application/vnd.oasis.opendocument.spreadsheet ,
  }

%% Embeds file #1 as #2 (the name a reader sees), with description #3 and relationship #4 (Data,
%% Supplement, Source, ...), as PDF object #5; lists it in the attachments panel.
\int_new:N \g__thesis_attach_int
\cs_new_protected:Npn \__thesis_attach:nnnnn #1#2#3#4#5
  {
    \group_begin:
      \pdfdict_put:nne { l_pdffile/Filespec } { AFRelationship } { /#4 }
      \pdf_string_from_unicode:nnN { utf8/string } {#3} \l_tmpa_str
      \pdfdict_put:nne { l_pdffile/Filespec } { Desc } { \l_tmpa_str }
      \pdffile_embed_file:nnn {#1} {#2} {#5}
    \group_end:
    \pdfmanagement_add:nne { Catalog / Names } { EmbeddedFiles } { \pdf_object_ref:n {#5} }
  }
\cs_generate_variant:Nn \__thesis_attach:nnnnn { eeeee }
%% The name a reader sees: the file name without its folders.
\cs_new_protected:Npn \__thesis_basename:nN #1#2
  {
    \file_parse_full_name:nNNN {#1} \l__thesis_dir_str \l__thesis_name_str \l__thesis_ext_str
    \tl_set:Ne #2 { \l__thesis_name_str \l__thesis_ext_str }
  }
\str_new:N \l__thesis_dir_str \str_new:N \l__thesis_name_str \str_new:N \l__thesis_ext_str
\tl_new:N \l__thesis_attachname_tl
%% Attached to the document as a whole (the catalogue's list of associated files).
\cs_new_protected:Npn \__thesis_attach_document:nnn #1#2#3 % file, description, relationship
  {
    \int_gincr:N \g__thesis_attach_int
    \__thesis_basename:nN {#1} \l__thesis_attachname_tl
    \__thesis_attach:eeeee {#1} { \l__thesis_attachname_tl } {#2} {#3}
      { thesis/attach/ \int_use:N \g__thesis_attach_int }
    \pdfmanagement_add:nne { Catalog } { AF }
      { \pdf_object_ref:e { thesis/attach/ \int_use:N \g__thesis_attach_int } }
  }

%% ---------------------------------------------------------------------------------------------
%% 3. \supplementfile[description]{file}
\NewDocumentCommand \supplementfile { O{} m }
  {
    \file_if_exist:nTF {#2}
      { \__thesis_attach_document:nnn {#2} { \tl_if_blank:nTF {#1} { Supplementary~file } {#1} } { Supplement } }
      { \msg_error:nnn { thesis } { attach-missing } {#2} }
  }
\msg_new:nnn { thesis } { attach-missing } { File~`#1'~to~attach~not~found. }

%% ---------------------------------------------------------------------------------------------
%% 2. \figuredata[description]{file}: in the tagged PDF an associated file of the latest image (the last
%% structure element tagged Figure); untagged, an associated file of the document.
\cs_new_protected:Npn \__thesis_lastfigure:N #1 % #1 := number of the latest Figure structure element (0: none)
  {
    \int_set:Nn #1 { \int_use:c { c@g__tag_struct_abs_int } }
    \bool_until_do:nn
      {
        \int_compare_p:nNn {#1} < { 2 }
        || \str_if_eq_p:ee { \prop_item:cn { g__tag_struct_ \int_use:N #1 _prop } { S } } { /Figure }
      }
      { \int_decr:N #1 }
    \int_compare:nNnT {#1} < { 2 } { \int_zero:N #1 }
  }
\int_new:N \l__thesis_figure_int
\NewDocumentCommand \figuredata { O{} m }
  {
    \file_if_exist:nTF {#2}
      {
        \tl_set:Ne \l_tmpb_tl { \tl_if_blank:nTF {#1} { Data~of~the~figure } { \exp_not:n {#1} } }
        \bool_set_false:N \l_tmpa_bool
        \cs_if_exist:NT \__tag_struct_add_AF:ee { \tag_if_active:T { \bool_set_true:N \l_tmpa_bool } }
        \bool_if:NTF \l_tmpa_bool
          {
            \__thesis_lastfigure:N \l__thesis_figure_int
            \int_compare:nNnTF { \l__thesis_figure_int } = { 0 }
              { \msg_error:nn { thesis } { figuredata-no-figure } }
              {
                \int_gincr:N \g__thesis_attach_int
                \__thesis_basename:nN {#2} \l__thesis_attachname_tl
                \__thesis_attach:eeeee {#2} { \l__thesis_attachname_tl } { \l_tmpb_tl } { Data }
                  { thesis/attach/ \int_use:N \g__thesis_attach_int }
                % (as tagpdf's own key AF does, for a structure element already begun)
                \__tag_struct_add_AF:ee { \int_use:N \l__thesis_figure_int }
                  { \pdf_object_ref:e { thesis/attach/ \int_use:N \g__thesis_attach_int } }
                \__tag_struct_prop_gput:nne { \int_use:N \l__thesis_figure_int } { AF }
                  { [ \tl_use:c { g__tag_struct_ \int_use:N \l__thesis_figure_int _AF_tl } ] }
              }
          }
          { \__thesis_attach_document:nVn {#2} \l_tmpb_tl { Data } }
      }
      { \msg_error:nnn { thesis } { attach-missing } {#2} }
  }
\cs_generate_variant:Nn \__thesis_attach_document:nnn { nVn }
\msg_new:nnn { thesis } { figuredata-no-figure }
  { \token_to_str:N \figuredata\ belongs~after~the~\token_to_str:N \includegraphics\ of~the~figure~it~describes. }

%% ---------------------------------------------------------------------------------------------
%% 1. The citation file, written at the start of every compile and attached (citation.bib, citation.ris).
\iow_new:N \g__thesis_cite_iow
\cs_new:Npn \__thesis_plain:n #1 { \text_purify:n {#1} }
%% One BibTeX field, `  name     = {value},', the = signs lined up (names padded to 9 characters)
\cs_new_protected:Npn \__thesis_bib_field:nn #1#2
  {
    \iow_now:Ne \g__thesis_cite_iow
      {
        \c_space_tl \c_space_tl #1
        \prg_replicate:nn { 9 - \str_count:n {#1} } { \c_space_tl }
        = ~ \c_left_brace_str #2 \c_right_brace_str ,
      }
  }
\cs_new_protected:Npn \__thesis_citation_write:
  {
    \tl_set:Ne \l_tmpa_tl { \__thesis_plain:n { \thesisauthorlastname } \thesisyear }
    \tl_remove_all:Nn \l_tmpa_tl { ~ }
    % BibTeX
    \iow_open:Nn \g__thesis_cite_iow { \c_sys_jobname_str -citation.bib }
    \iow_now:Ne \g__thesis_cite_iow { @phdthesis \c_left_brace_str \l_tmpa_tl , }
    \__thesis_bib_field:nn { author } { \__thesis_plain:n { \thesisauthorlastname , ~ \thesisauthorfirstnames } }
    \__thesis_bib_field:nn { title } { \c_left_brace_str \__thesis_plain:n { \thesispdftitle } \c_right_brace_str }
    \__thesis_bib_field:nn { school } { Delft~University~of~Technology }
    \__thesis_bib_field:nn { address } { Delft,~the~Netherlands }
    \__thesis_bib_field:nn { type } { Doctoral~dissertation }
    \__thesis_bib_field:nn { year } { \thesisyear }
    \thesis_if_filled:nT { \thesisisbn } { \__thesis_bib_field:nn { isbn } { \thesisisbn } }
    \thesis_if_filled:nT { \thesisdoi } { \__thesis_bib_field:nn { doi } { \thesisdoi } }
    \__thesis_bib_field:nn { keywords } { \__thesis_plain:n { \thesiskeywords } }
    \iow_now:Ne \g__thesis_cite_iow { \c_right_brace_str }
    \iow_close:N \g__thesis_cite_iow
    % RIS
    \iow_open:Nn \g__thesis_cite_iow { \c_sys_jobname_str -citation.ris }
    \iow_now:Ne \g__thesis_cite_iow { TY\c_space_tl \c_space_tl -~THES }
    \iow_now:Ne \g__thesis_cite_iow { AU\c_space_tl \c_space_tl -~\__thesis_plain:n { \thesisauthorlastname , ~ \thesisauthorfirstnames } }
    \iow_now:Ne \g__thesis_cite_iow { TI\c_space_tl \c_space_tl -~\__thesis_plain:n { \thesispdftitle } }
    \iow_now:Ne \g__thesis_cite_iow { PY\c_space_tl \c_space_tl -~\thesisyear }
    \iow_now:Ne \g__thesis_cite_iow { PB\c_space_tl \c_space_tl -~Delft~University~of~Technology }
    \iow_now:Ne \g__thesis_cite_iow { CY\c_space_tl \c_space_tl -~Delft,~the~Netherlands }
    \iow_now:Ne \g__thesis_cite_iow { M3\c_space_tl \c_space_tl -~Doctoral~dissertation }
    \thesis_if_filled:nT { \thesisisbn } { \iow_now:Ne \g__thesis_cite_iow { SN\c_space_tl \c_space_tl -~\thesisisbn } }
    \thesis_if_filled:nT { \thesisdoi }
      {
        \iow_now:Ne \g__thesis_cite_iow { DO\c_space_tl \c_space_tl -~\thesisdoi }
        \iow_now:Ne \g__thesis_cite_iow { UR\c_space_tl \c_space_tl -~https://doi.org/\thesisdoi }
      }
    \clist_map_inline:en { \__thesis_plain:n { \thesiskeywords } }
      { \iow_now:Ne \g__thesis_cite_iow { KW\c_space_tl \c_space_tl -~##1 } }
    \iow_now:Ne \g__thesis_cite_iow { ER\c_space_tl \c_space_tl -~ }
    \iow_close:N \g__thesis_cite_iow
  }
\cs_generate_variant:Nn \clist_map_inline:nn { en }
\AddToHook{begindocument/end}
  {
    \__thesis_citation_write:
    \int_gincr:N \g__thesis_attach_int
    \__thesis_attach:eeeee { \c_sys_jobname_str -citation.bib } { citation.bib }
      { How~to~cite~this~dissertation~(BibTeX) } { Supplement } { thesis/attach/ \int_use:N \g__thesis_attach_int }
    \pdfmanagement_add:nne { Catalog } { AF } { \pdf_object_ref:e { thesis/attach/ \int_use:N \g__thesis_attach_int } }
    \int_gincr:N \g__thesis_attach_int
    \__thesis_attach:eeeee { \c_sys_jobname_str -citation.ris } { citation.ris }
      { How~to~cite~this~dissertation~(RIS:~EndNote,~Zotero,~Mendeley) } { Supplement }
      { thesis/attach/ \int_use:N \g__thesis_attach_int }
    \pdfmanagement_add:nne { Catalog } { AF } { \pdf_object_ref:e { thesis/attach/ \int_use:N \g__thesis_attach_int } }
  }
\ExplSyntaxOff
 % files attached to the PDF: citation file, figure data, supplements

%% The template's own commands (template/commands.tex): \repeatcaption, \file, \doi, \orcid,
%% \chapterbasedon, \panelref, the list layout and the spacing of \inf, \max, ...
%% =====================================================================
%% The template's own commands (read by template/packages.tex). Your own commands and notation go in
%% settings/commands.tex.
%% =====================================================================

%% The same figure a second time, with its original number: \repeatcaption{fig:label}{caption}
%% (the caption without a full stop: "(repeated from page 13)." follows it)
\newcommand{\repeatcaption}[2]{%
  \renewcommand{\thefigure}{\ref{#1}}%
  \renewcommand{\theHfigure}{repeated.\arabic{chapter}.\arabic{figure}}% (a link target of its own: the next figure gets this number)
  \renewcommand{\addcontentsline}[3]{}% not in the list of figures a second time
  \caption{#2 (repeated from page~\pageref{#1}).}% caption without a full stop: one is added
  \addtocounter{figure}{-1}% So that next figure after the repeat gets the right number.
}

%% inf, max, min, lim and lim inf get the depth of a descender, as sup and lim sup have from their p, so
%% the limits under operators side by side line up: \inf_{q} \sup_{v} puts q and v on one baseline.
\AtBeginDocument{%
    \renewcommand{\inf}{\operatorname*{inf\vphantom{p}}}%
    \renewcommand{\max}{\operatorname*{max\vphantom{p}}}%
    \renewcommand{\min}{\operatorname*{min\vphantom{p}}}%
    \renewcommand{\lim}{\operatorname*{lim\vphantom{p}}}%
    \renewcommand{\liminf}{\operatorname*{lim\,inf\vphantom{p}}}}

%% Lists: dashes as bullets, no extra space between items (a list reads as part of the paragraph),
%% and the item text at the paragraph indent, with the dash or number hanging in front of it.
\tightlists
\renewcommand{\labelitemi}{\textendash}\setlength{\leftmargini}{\parindent}
\renewcommand{\labelitemii}{\textendash}\setlength{\leftmarginii}{\parindent}
\renewcommand{\labelitemiii}{\textendash}\setlength{\leftmarginiii}{\parindent}
\renewcommand{\labelitemiv}{\textendash}\setlength{\leftmarginiv}{\parindent}

%% Dutch name prefixes: write {\VAN} in references.bib (e.g. author = {T. \VAN~Pelt}). It prints "Van"
%% in the text (\citet) and "van" in the reference list (switched in main.tex).
\DeclareRobustCommand{\VAN}[0]{Van}

%% File and package names in typewriter type that may break across lines after . / _ (a long name in
%% \texttt can't be broken and sticks out into the margin on a narrow page). Type underscores as they are:
%% \file{settings/nomenclature.tex}, \file{figures/plots/thesis.mplstyle}
\newcommand{\file}[1]{\nolinkurl{#1}}
%% Identifiers outside the references, in typewriter and as links, set like the DOIs in the references:
%% \doi{10.4121/00000000} gives "DOI: 10.4121/00000000", \orcid{0000-0002-1825-0097} the ORCID icon
%% and the iD, both linking to the ORCID record.
\newcommand{\doi}[1]{\textsc{doi}:\nobreakspace\href{https://doi.org/#1}{\expandafter\nolinkurl\expandafter{#1}}}
\newcommand{\orcid}[1]{\expandafter\orcidlink\expandafter{#1}\,\href{https://orcid.org/#1}{\expandafter\nolinkurl\expandafter{#1}}}

%% Chapter based on a paper: a small note at the start of the chapter text with the reference and the
%% author contributions. \chapterbasedon{\cite{key}}{A.~Author: conceptualisation, ...; B.~Author: ...}
\makeatletter
\newcommand{\chapterbasedon}[2]{%
    \par\begingroup\small\setlength{\parindent}{0pt}%
    This chapter is based on #1.\par
    \textit{Author contributions:} #2\par
    \endgroup\bigskip
    \@afterindentfalse\@afterheading % the text after the note starts unindented
}
\makeatother

%% A panel of a figure as one link, letter included: \panelref{fig:x}{a} gives "Figure 2.2a"
%% (\Panelref at the start of a sentence; several panels: \panelref{fig:x}{a,b}).
\newcommand{\panelref}[2]{\hyperref[#1]{\cref*{#1}#2}}
\newcommand{\Panelref}[2]{\hyperref[#1]{\Cref*{#1}#2}}


%% The parts of the book: \frontmatter, \mainmatter, \printreferences, \appendix, \backmatter
%% (template/matter.tex)
%% =====================================================================
%% The parts of the book, as main.tex lists them (read by template/packages.tex):
%%   \frontmatter    page numbers in roman from the half-title, the title pages (template/titlepage.tex,
%%                   from settings/thesisinfo.tex), then the summary, preface, contents, nomenclature:
%%                   no running heads, a compact block in the contents
%%   \tableofcontents, \listoffigures, \listoftables, \listofalgorithms
%%                   each on a new page; the contents get a bookmark but no entry in themselves
%%   \mainmatter     page numbers in arabic from 1, chapters with running heads and thumb tabs
%%   \printreferences  the reference list at the end of the book (nothing with the chapterbib option:
%%                   every chapter then ends with its own), without thumb tabs
%%   \appendix       chapters lettered A, B, ...: "Appendix A" in references, bookmarks at the top level,
%%                   thumb tabs continuing below the last chapter
%%   \backmatter     acknowledgements, CV, publications: as the front matter
%% The line spacing of the text (\thesisspacing, set per page size in thesis.cls) starts after the title
%% pages.
%% =====================================================================
\makeatletter
\setcounter{tocdepth}{1} % the contents go down to sections (\setcounter{tocdepth}{2} in main.tex: subsections too)

\newif\ifthesis@spacing
%% The book ends: the last page out, and the line spacing of the text ends with it (before the digital
%% edition's back cover, template/digitalcover.tex)
\AddToHook{enddocument}[thesis/matter]{\ifthesis@spacing\clearpage\par\endgroup\fi}
\DeclareHookRule{enddocument}{thesis/matter}{before}{thesis/backcover}
\renewcommand\frontmatter{%
    \pagenumbering{roman}% counted from the half-title (page i), as in books; the title pages print no number
    %% Title pages of a TU Delft dissertation. All names, dates and numbers come from settings/thesisinfo.tex.
%%   page i   : half-title
%%   page ii  : blank
%%   page iii : official title page (wording prescribed by TU Delft; check it against the
%%              Graduate School's current template before printing)
%%   page iv  : doctoral committee and colophon
%%   page v   : dedication (optional; \thesisdedication in settings/thesisinfo.tex)
\makeatletter
%% Centred text starting at the height of the chapter titles
\newcommand{\thesis@titleblock}[1]{\begingroup \parskip\z@ \parindent\z@
    \chapterheadstart\thesis@chapterlabelspace \centering #1\endgroup}

%% ---------------------------------------------------------------------
%% Page i: half-title: the title alone (the subtitle waits for the title page), smaller than there
%% ---------------------------------------------------------------------
\thispagestyle{empty}
\thesis@titleblock{\fontsize{0.75\dimexpr\chaptertitlesize\relax}{1.15\dimexpr\chaptertitlesize\relax}\headingfont
    \thesistitle\par}

%% ---------------------------------------------------------------------
%% Page ii: blank
%% ---------------------------------------------------------------------
\cleardoublepage

%% ---------------------------------------------------------------------
%% Page iii: official title page
%% ---------------------------------------------------------------------
\thispagestyle{empty}
\thesis@titleblock{%
    {\fontsize{\chaptertitlesize}{1.15\dimexpr\chaptertitlesize\relax}\headingfont\thesistitle\par}
    \bigskip
    {\sffamily\Large\thesissubtitle\par}}
\vfill
\begin{center}
    {\runningheadcaps DISSERTATION}

    \bigskip

    for the purpose of obtaining the degree of doctor \\
    at Delft University of Technology \\
    by the authority of the Rector Magnificus, \rectormagnificus, \\
    chair of the Board for Doctorates \\
    to be defended publicly on \\
    \thesisdefencedate

    \bigskip
    \bigskip

    by

    \bigskip
    \bigskip

    {\headingfont\Large\bfseries \thesisauthorfirstnames\ \MakeUppercase{\thesisauthorlastname}}

    \bigskip
    \bigskip

    \thesisdegree \\
    born in \thesisbirthplace

    \vspace*{2\bigskipamount}
\end{center}

%% ---------------------------------------------------------------------
%% Page iv: doctoral committee and colophon
%% ---------------------------------------------------------------------
\clearpage
\thispagestyle{empty}
{\small
\setlength{\parskip}{0pt}
\noindent This dissertation has been approved by the promotors.

\medskip
\noindent Composition of the doctoral committee:

\smallskip
\noindent
\begin{tabular}{@{}p{0.5\textwidth}@{}p{0.5\textwidth}@{}}
    Rector Magnificus,  & chairperson \\
    \thesiscommittee
\end{tabular}

\IfThesisEmpty{\thesiscommitteenote}{}{%
    \medskip
    \noindent \thesiscommitteenote}

\vfill
\noindent\includegraphics[height=8mm, alt={Delft University of Technology logo}]{template/images/tudelft-logo.pdf}

\bigskip
\IfThesisEmpty{\thesisfunding}{}{%
    \noindent \thesisfunding\par\medskip}
\noindent
\begin{tabular}{@{}p{0.22\textwidth}@{}p{0.78\textwidth}@{}}
    \textit{Keywords:}   & \thesiskeywords \\[\smallskipamount]
    \IfThesisEmpty{\thesisdata}{}{\textit{Data:} & \thesisdata \\[\smallskipamount]}
    \IfThesisEmpty{\thesiscode}{}{\textit{Code:} & \thesiscode \\[\smallskipamount]}
    \textit{Printed by:} & \thesisprinter  \\[\smallskipamount]
    \textit{Cover by:}   & \thesiscover    \\
\end{tabular}

\bigskip
\ifdefstring{\thesislicence}{CC0 1.0}{% public domain: no copyright line
\noindent This work is dedicated to the public domain under \href{\thesislicenceurl}{CC0~1.0~Universal}.
}{%
\noindent Copyright \textcopyright\ \thesisyear\ by \thesisauthorinitials~\thesisauthorlastname
\par{\raggedright\noindent % (licence names are long words: Attribution-NonCommercial-ShareAlike)
\IfThesisEmpty{\thesislicence}{All rights reserved.}{This work is licensed under a \thesislicencename\ licence
    (\href{\thesislicenceurl}{\thesislicencecode}).}\par}}

\smallskip
\noindent ISBN \thesisisbn

\smallskip
\noindent An electronic copy of this dissertation is available at \\
\IfThesisEmpty{\thesisdoi}{\url{https://repository.tudelft.nl}}% (links to the dissertation's DOI once known)
    {\href{https://doi.org/\thesisdoi}{\nolinkurl{https://repository.tudelft.nl}}}.
}
\cleardoublepage

%% ---------------------------------------------------------------------
%% Page v (optional): dedication, at the height of the half-title, with a blank page behind it
%% ---------------------------------------------------------------------
\IfThesisEmpty{\thesisdedication}{}{%
    \thispagestyle{empty}%
    \thesis@titleblock{\itshape\thesisdedication\par}%
    \cleardoublepage}
\makeatother
%
    \par\begingroup\def\baselinestretch{\thesisspacing}\@currsize % the line spacing of the text, up to the end
    \global\thesis@spacingtrue
    \pagestyle{frontmatter}%
    \aliaspagestyle{chapter}{empty}%
    \contentsstyle{frontmatter}%
    \global\thesis@starredtoctrue} % \chapter*: in the contents (template/formatText.tex)

%% \tableofcontents: on a new page, with a bookmark of its own (it does not list itself).
%% (\tableofcontents* is the bare list, as main.tex had it before: the page break and bookmark were there.)
\DeclareDocumentCommand\tableofcontents{s}{%
    \IfBooleanF{#1}{\clearpage\pdfbookmark[0]{\contentsname}{contents}}%
    \thesis@list{\BooleanTrue}{toc}{\contentsname}{\contentsname}}
\DeclareDocumentCommand\listoffigures{s}{\clearpage\thesis@list{#1}{lof}{\listfigurename}{\listfigurename}}
\DeclareDocumentCommand\listoftables{s}{\clearpage\thesis@list{#1}{lot}{\listtablename}{\listtablename}}
\AddToHook{cmd/listofalgorithms/before}{\clearpage}

\renewcommand\mainmatter{%
    \pagenumbering{arabic}%
    \aliaspagestyle{chapter}{chapteropening}%
    \contentsstyle{mainmatter}%
    \global\thesis@starredtocfalse}

\NewDocumentCommand\printreferences{}{%
    \ifchapterbib\else
        \thumbsoff % no thumb tab on the reference pages
        \aliaspagestyle{chapter}{empty}%
        \DeclareRobustCommand{\VAN}[0]{van}% "van" prefixes lowercase in the reference list (template/commands.tex)
        \renewcommand{\bibname}{References}%
        \printbibliography
        \pagestyle{chapteronly}%
    \fi
    \DeclareRobustCommand{\VAN}[0]{Van}}% and capitalised again in the text

\AddToHook{cmd/appendix/before}{%
    \bookmarksetup{startatroot}% the appendices' bookmarks at the top level, not under References
    \aliaspagestyle{chapter}{chapteropening}}
\AddToHook{cmd/appendix/after}{%
    \crefalias{chapter}{appendix}\crefalias{section}{subappendix}% \cref gives "Appendix A", "Section A.1"
    \thumbson % appendices get tabs, continuing below the last chapter
    \global\thesis@inappendixtrue % their page style: appendix (template/formatText.tex)
    \clearpage}

\renewcommand\backmatter{%
    \bookmarksetup{startatroot}%
    \cleardoublepage
    \aliaspagestyle{chapter}{empty}%
    \pagestyle{frontmatter}%
    \contentsstyle{frontmatter}%
    \global\thesis@starredtoctrue}
\makeatother

